%%
%% Copyright 2007-2024 Elsevier Ltd
%%
%% This file is part of the 'Elsarticle Bundle'.
%% ---------------------------------------------
%%
%% It may be distributed under the conditions of the LaTeX Project Public
%% License, either version 1.3 of this license or (at your option) any
%% later version.  The latest version of this license is in
%%    http://www.latex-project.org/lppl.txt
%% and version 1.3 or later is part of all distributions of LaTeX
%% version 1999/12/01 or later.
%%
%% The list of all files belonging to the 'Elsarticle Bundle' is
%% given in the file `manifest.txt'.
%%
%% Template article for Elsevier's document class `elsarticle'
%% with numbered style bibliographic references
%% SP 2008/03/01
%% $Id: elsarticle-template-num.tex 249 2024-04-06 10:51:24Z rishi $
%%
\documentclass[preprint,12pt]{elsarticle}

%% Use the option review to obtain double line spacing
%% \documentclass[authoryear,preprint,review,12pt]{elsarticle}

%% Use the options 1p,twocolumn; 3p; 3p,twocolumn; 5p; or 5p,twocolumn
%% for a journal layout:
%% \documentclass[final,1p,times]{elsarticle}
%% \documentclass[final,1p,times,twocolumn]{elsarticle}
%% \documentclass[final,3p,times]{elsarticle}
%% \documentclass[final,3p,times,twocolumn]{elsarticle}
%% \documentclass[final,5p,times]{elsarticle}
%% \documentclass[final,5p,times,twocolumn]{elsarticle}

%% For including figures, graphicx.sty has been loaded in
%% elsarticle.cls. If you prefer to use the old commands
%% please give \usepackage{epsfig}
%% The amssymb package provides various useful mathematical symbols
\usepackage{amssymb}
%% The amsmath package provides various useful equation environments.
\usepackage{amsmath,amsthm,proof,float}
\usepackage{semantic}
\usepackage{mathtools}
\usepackage{listings}
\usepackage{hyperref}
\def\lstlanguagefiles{lstlean.tex}
% set default language
\lstset{language=lean}
\newtheorem{Lemma}{Lemma}
\newtheorem{Theorem}{Theorem}
\theoremstyle{definition}
\newtheorem{Example}{Example}

%% Proof sketches: same layout as amsthm's proof environment, but headed by
%% "Proof sketch" and closed by a filled square, since the argument is informal.
\newenvironment{proofsketch}
  {\renewcommand{\qedsymbol}{$\blacksquare$}\begin{proof}[Proof sketch]}
  {\end{proof}}

%% The amsthm package provides extended theorem environments
%% \usepackage{amsthm}

%\usepackage{xcolor}

\usepackage{color}
\definecolor{keywordcolor}{rgb}{0.7, 0.1, 0.1}   % red
\definecolor{tacticcolor}{rgb}{0.0, 0.1, 0.6}    % blue
\definecolor{commentcolor}{rgb}{0.4, 0.4, 0.4}   % grey
\definecolor{symbolcolor}{rgb}{0.0, 0.1, 0.6}    % blue
\definecolor{sortcolor}{rgb}{0.1, 0.5, 0.1}      % green
\definecolor{attributecolor}{rgb}{0.7, 0.1, 0.1} % red

\def\lstlanguagefiles{lstlean.tex}
% set default language
\lstset{language=lean}

% This is a literate Agda document (see
% https://agda.readthedocs.io/en/latest/tools/literate-programming.html).
% Build with:
%   agda -i . -i ../../src/pest-sn/src --latex paper.lagda.tex
%   xelatex latex/paper.tex   (a Unicode-aware engine is required)
% Agda copies agda.sty into the generated latex/ directory.
\usepackage{agda}

% Unicode-capable monospace font for the Agda code blocks (LuaLaTeX/XeLaTeX).
\usepackage{iftex}
\ifPDFTeX\else
  \usepackage{fontspec}
  \setmonofont{DejaVu Sans Mono}[Scale=MatchLowercase]
  % agda.sty renders symbols/keywords in sans and bound variables in italic,
  % neither of which covers the Unicode we use; route every token through the
  % (Unicode) monospace font instead.
  \renewcommand{\AgdaFontStyle}[1]{\texttt{#1}}
  \renewcommand{\AgdaKeywordFontStyle}[1]{\texttt{#1}}
  \renewcommand{\AgdaBoundFontStyle}[1]{\texttt{#1}}
  % Fallback for the few glyphs DejaVu Sans Mono lacks (e.g. the script R 𝓡).
  % Loaded by file name: the system font "DejaVu Math TeX Gyre" has no bold or
  % italic cut, and asking fontspec for it by family name makes LuaTeX abort
  % while embedding the substituted shapes.
  \usepackage{newunicodechar}
  \newfontfamily\AgdaFallback{texgyredejavu-math.otf}
  \newunicodechar{𝓡}{{\AgdaFallback 𝓡}}
\fi

\usepackage{lineno}
\usepackage{url}

%% The lineno packages adds line numbers. Start line numbering with
%% \begin{linenumbers}, end it with \end{linenumbers}. Or switch it on
%% for the whole article with \linenumbers.
%% \usepackage{lineno}

\journal{Journal of Computer Languages}

\begin{document}

\begin{frontmatter}

%% Title, authors and addresses

%% use the tnoteref command within \title for footnotes;
%% use the tnotetext command for theassociated footnote;
%% use the fnref command within \author or \affiliation for footnotes;
%% use the fntext command for theassociated footnote;
%% use the corref command within \author for corresponding author footnotes;
%% use the cortext command for theassociated footnote;
%% use the ead command for the email address,
%% and the form \ead[url] for the home page:
%% \title{Title\tnoteref{label1}}
%% \tnotetext[label1]{}
%% \author{Name\corref{cor1}\fnref{label2}}
%% \ead{email address}
%% \ead[url]{home page}
%% \fntext[label2]{}
%% \cortext[cor1]{}
%% \affiliation{organization={},
%%             addressline={},
%%             city={},
%%             postcode={},
%%             state={},
%%             country={}}
%% \fntext[label3]{}

\title{Pest control: A formal model of the Pest parser generator}
%% use optional labels to link authors explicitly to addresses:
%% \author[label1,label2]{}
%% \affiliation[label1]{organization={},
%%             addressline={},
%%             city={},
%%             postcode={},
%%             state={},
%%             country={}}
%%
%% \affiliation[label2]{organization={},
%%             addressline={},
%%             city={},
%%             postcode={},
%%             state={},
%%             country={}}

\author[ufop]{Guilherme Daher} %% Author name
\author[ufop]{Elton Cardoso}
\author[ufsc]{Samuel Feitosa}
\author[ufop]{Rodrigo Ribeiro} %% Author name

%% Author affiliation
\affiliation[ufop]{organization={Federal University of Ouro Preto},%Department and Organization
            city={Ouro Preto},
            state={MG},
            country={Brazil}}
\affiliation[ufsc]{organization={Federal University of the Southern Border},%Department and Organization
            city={Chapecó},
            state={SC},
            country={Brazil}}

%% Abstract
\begin{abstract}
Parsing expressions grammars (PEGs) are a recognition-based formalism for parsers,
which has been gaining popularity because they avoid some problems present in context-free grammar
based specifications. In the context of the Rust programming language, the most used
PEG-based parser generator is Pest. An interesting feature of Pest grammars is that they
support new operators for repetition and handling an implicit stack during parsing.
In this work, we propose a formalization of these new constructs and extend a
type inference algorithm which guarantees that all well-typed Pest-style grammars
terminate on all inputs. We prove this termination guarantee as a machine-checked
\emph{strong normalization} theorem in the Agda proof assistant, using a
logical-relations argument, and present the development as a literate Agda
rendering of this article.
\end{abstract}

%%Graphical abstract
%\begin{graphicalabstract}
%\includegraphics{grabs}
%\end{graphicalabstract}

%%Research highlights
\begin{highlights}
\item Formalizes the semantics of Pest parser generator new repetition and stack operators
  using an operational semantics.
\item Extends an existing PEG type system and its inference algorithm that ensures
  termination for all well-typed grammars to consider Pest new operators.
\item We provide an Agda mechanized proof of \emph{strong normalization} for the defined type
  system: parsing every well-typed grammar terminates on every input, established by a
  logical-relations argument with no appeal to a fuel counter.
\end{highlights}

%% Keywords
\begin{keyword}
Parsing \sep parsing expression grammars \sep type systems
\end{keyword}

\end{frontmatter}

%% Add \usepackage{lineno} before \begin{document} and uncomment
%% following line to enable line numbers.
%% NOTE: disabled because the lineno package conflicts with the literate Agda
%% code environments (\@LN@do@vadjusts vs. \@vspace@calcify). Re-enable only
%% for a non-literate build.
% \linenumbers

%% main text
%%

% Literate Agda module header and imports of the mechanised development.
% Hidden from the rendered PDF; typechecked against the pest-sn library.
\begin{code}[hide]
module paper where

open import Data.List using ([])
open import Data.Fin.Subset using (_∉_)
open import Data.Product using (∃; _,_)
\end{code}

\section{Introduction}\label{sec:intro}

Parsing is a core problem in computer science. In simple terms, a parsing
algorithm, also named a parser, checks if a string of symbols conforms to some
set of rules specified by formalisms such as context-free grammars (CFGs) or
regular expressions (REs). Usually, parsers construct data structures showing
what rules were used to conclude that the input string was obtainable from the
specification or an error message indicating that such string is unobtainable
from the specification. Due to its importance, parsing is an extensively studied
subject, as evidenced by Grune and Jacob's book "Parsing Techniques: A Practical
Guide", which has more than 600 pages and covers only classic algorithms \cite{Grune90}.
However, parsing is far from being a solved problem: in recent years,
we have seen a growing interest in the development of new algorithms and
formalisms~\cite{Ford2004, Trevor10,Asperti10, Lucks17, Zhang23} and the verification
of well-known parsing techniques~\cite{FirsovU14, FirsovU13,LopesRC16,Bernardy16,JourdanPL12}.

In the early 2000's, Brian Ford proposed \emph{Parsing Expression Grammars}
(PEGs),  a recognition-based formalism to
define languages~\cite{Ford2004}. The novelty of PEGs is that they
specify how to check whether a string is in
the language being defined and not how it is generated by its specification.
Formally a PEG is a 4-tuple $(V, \Sigma, R, e_S)$ where $V$ is the set of
non-terminals (also called variables), $\Sigma$ is the alphabet of terminals,
$R$ is a function which maps a non terminal from $V$ to a {\it parsing expression}
and $e_S$ is the initial parsing
expression. We use this same notation $(V, \Sigma, R, e_S)$ throughout the paper. A parsing expression is very
similar to the right-hand side of a CFG in the extended Backus-Naur
form with the addition of new operators such as the not-predicate (!) and
the prioritized choice (/). A fundamental difference between PEG
and CFGs is the choice operator, which imposes priority between alternatives.
It means that $\beta$ in a prioritized choice expression $\alpha / \beta$ is tested only
in case of failure on $\alpha$. The not-predicate operator checks
if the string matches some syntax without consuming the input.

In the context of the Rust programming language~\cite{RustBook}, the
Pest library~\cite{PEST} is one of the most popular PEG-based parser
generators for Rust, with more than 200 million downloads registered at
the \href{http://crates.io}{crates.io} package registry as of early
2026\footnote{Retrieved from \url{https://crates.io/crates/pest} in early
2026.}. Besides supporting
all traditional PEG operators, Pest supports operators to manipulate an
implicit stack during parser execution~\cite{PESTSTACK}. The inclusion of a stack allows
Pest parsers to elegantly express indentation sensitive languages by
controlling indentation levels using the stack. While such new operators
ease the task of writing grammars for complex formats, additional care
must be taken in order to ensure PEG \emph{well-formedness}. We
say that a PEG is well-formed if it terminates on any input
string~\cite{Ford2004}. Pest users have reported that the tool
fails to detect grammars that loop on some inputs~\cite{PESTBUG}. In the
discussion thread that reports this issue, the author conjectures that detecting
such looping grammars might be undecidable in general.
In this work, we provide a sound (if necessarily incomplete)
static criterion that rejects the looping grammars while accepting a large and
useful class. Our approach consists of formalizing a generalized
version of these Pest operators using operational semantics and
an extended version of a type system for PEGs~\cite{Ribeiro19} and its
inference algorithm~\cite{Cardoso23} to deal with these operators.

Specifically, we make the following contributions:

\begin{itemize}
  \item We formalize both stack and the new repetitions Pest operators using
    operational semantics and an extended version of a type system for PEGs~\cite{Ribeiro19}
    to deal with these new operations. We prove, in the Agda proof
    assistant~\cite{agda}, a \emph{strong normalization} theorem for well-typed grammars (parsing terminates on every input), using a logical-relations argument in the style
    of Tait's proof for the simply-typed $\lambda$-calculus~\cite{plfa22.08}, so that
    termination is a machine-checked theorem rather than an informal consequence of a
    fuel-based interpreter. The complete formalization source code is available
    on-line~\cite{pegrepository}.
  \item We include operations to query the top-most element of the stack
    which can be used to handle context-sensitive formats, such as length
    indexed fields, common in binary data formats;
  \item We show how to extend a type inference algorithm~\cite{Cardoso23} for PEGs
    to deal with Pest's new operators. A prototype implementation of the inference
    algorithm was developed using the Racket programming language~\cite{Culpepper19}
    and used to check some example grammars.
\end{itemize}

This paper is an extended version of our conference paper~\cite{Daher25},
presented at the 2025 Brazilian Symposium on Programming Languages (SBLP 2025).
For the benefit of the reader, we state explicitly which material is reproduced
from the conference version and which is new. The overview of Pest operators
(Section~\ref{sec:pest-operators}), the review of PEG syntax, semantics and the
type system for PEGs (the corresponding parts of Section~\ref{sec:background}),
and the operational semantics and typing rules for the new operators
(Section~\ref{sec:pest-formalization}, excluding the mechanization) are, up to
minor revisions, reproduced from~\cite{Daher25}. The genuinely new contributions
of this paper are: 1) the complete formalization of the type system and the
relational semantics in Agda; 2) a machine-checked
strong-normalization theorem for Pest-style grammars, covering the stack, exact
and interval repetition operators; 3) a
larger set of examples (Section~\ref{sec:examples}); and 4) an empirical
evaluation of the inference algorithm on a corpus of real Pest grammars.

The rest of this paper is organized as follows.
Section~\ref{sec:pest-operators} presents an informal description of Pest operators.
Section~\ref{sec:background}
reviews the syntax, semantics and a type system for PEGs.
Section~\ref{sec:pest-formalization} defines the syntax, semantics, the type
system and its inference algorithm for the new PEG operators. Section~\ref{sec:typesoundnesslean}
describes our mechanised strong-normalization proof for Pest grammars in Agda.
Section~\ref{sec:examples}
presents some examples of grammars accepted by our tool.
Related work is discussed on Section~\ref{sec:related} and
Section~\ref{sec:conclusion} concludes.


\section{An overview of Pest new operators}~\label{sec:pest-operators}

Pest is a mature parser generator for the Rust ecosystem: a grammar is written in
a dedicated file using a PEG-like syntax, from which Pest derives a parser. Two
features distinguish it from a textbook PEG. First, Pest provides an \emph{implicit
stack} as part of its core functionality (it is a built-in feature of the tool,
not an extension introduced by us), together with operators that push, pop and
peek at it; the stack is what lets Pest describe context-sensitive formats such as
indentation-sensitive or length-prefixed languages. Second, Pest offers repetition
operators with explicit bounds. In this paper we formalise a generalised version
of these two families of operators. For the stack, we additionally propose two
indexes, \textbf{top.tonat} and \textbf{top.length}, that inspect the top-of-stack
element; these two indexes are our own additions and are not currently part of
Pest. We now give an informal account of the operators before their formal
treatment in Section~\ref{sec:pest-formalization}.

Before starting the description of each new operator, we describe
\emph{repetition indexes}, which are values ranging over the natural numbers and a
special value, \textbf{infty}, which denotes an upper-bound over
indexes\footnote{For any index $i \neq \mathbf{infty}$, we have
that $i < \mathbf{infty}$ and $i \circ \mathbf{infty} = \infty$}.
The notation $\circ$ denotes either addition
or multiplication operations over natural numbers. The constants
\textbf{top.tonat} and \textbf{top.length} refer to the conversion
of the stack top element to a numeric value and its size, respectively.
More details on how these operations work will be described
on the semantics.

The first operator we consider is the iterated repetition, denoted $e^i$ in
the abstract syntax and \verb|e^i| in the concrete syntax. It describes that
expression $e$ should be repeated $i$ times over the input, failing if any of these iterations
results in a parser failure. As a first example, consider the task of specifying a format formed by two byte
values followed by a single bit terminator. We could express this format with the
following grammar using the iterated repetition:
\begin{verbatim}
bit <-- '0' / '1'
byte <-- bit^8
start: byte^2 ~ bit
\end{verbatim}
The notation \verb|'0'| denotes an expression for a string,
operator \verb|~| denotes concatenation (following Pest concrete syntax),
\verb|A <-- e| denotes a grammar rule and \verb|start:| specifies
the starting expression for this grammar.
For such small counts the operator is mainly a convenience, since \verb|byte^2|
could equally be written \verb|byte ~ byte|. Its real value appears with the
large, fixed field widths that pervade binary and network formats, where writing
the sub-expression out explicitly would be verbose and error-prone. As example,
consider the task of writing an expression for a 16-byte IPv6 address or 
a 32-byte SHA-256 digest, which could be simply denoted by \verb|byte^16| and
\verb|byte^32|. In these cases the iterated repetition
operator states the exact width once, mirroring the way such fields are
described in their specifications, which would be impractical to spell out by hand.

The Pest tool also introduces an interval repetition operator, $e[i,j]$,
which executes the expression $e$ a number of times between the indexes $i$
and $j$, inclusive. As an example of this operator usage, consider the
task of parsing either 2, 3 or 4 bits sequences from input. This could be
expressed elegantly by the following expression using the interval repetition:
\verb|('0' / '1')[2,4]|.

However, not all combinations of indexes are allowed in repetition operators.
An example of an invalid repetition operator: \verb|e[infty, 2]|, since a
bounded repetition, \verb|e[i,j]|, is valid if $i < j$ and $i$ is not the constant
\verb|infty|. Regarding the exact repetition operator, $e^i$, index $i$ is not allowed to
be \textbf{infty}, since its semantics is to parse the input by repeating
an expression $e$ \textbf{exactly} $i$ times, which makes the expression
$e^{\mathbf{infty}}$ to be repeat indefinitely over the input.

Now, we turn our attention to the stack manipulation operators. Expression
\verb|push(e)| parses the input using the expression \texttt{e} and pushes the
consumed input on the parser stack. Operator \verb|pop| removes the current
stack top-most element and tries to match it on the current input,
\verb|peek| also matches the input without removing the element from the
stack and \verb|drop| removes the first element of current stack, without
matching it. The operator \verb|peekall| tries to match the input using the
complete stack content, without emptying the stack. Finally, \verb|dropall|
empties the stack.

One interesting example use of stack operators is to parse
length-indexed fields present in several data formats.
The next grammar shows how to express netstrings, which are a format method
for strings that includes a field denoting the size of the represented data.
\begin{verbatim}
number <-- '0' / '1' / ... / '9'
letter <-- 'a' / ... / 'z'
char <-- number / letter

start: push(number+) ~ ':' ~ char^top.tonat ~ ','
\end{verbatim}
The starting expression begins by pushing the input size value on the parser stack,
and then access this size information using \verb|top.tonat| to parse the exact
quantity of input based on the size information that starts the netstring.
Operators \verb|top.tonat| and \verb|top.length| are not present in the current
definition of Pest parser generator. We include both since they have a direct
formalization and increases the expressivity for Pest grammars which can be used
to represent length-indexed fields, which are notoriously difficult to parse
correctly~\cite{Zhang23, Lucks17}.


\section{Background}\label{sec:background}

\subsection{An Overview of PEGs}

Intuitively, PEGs are a formalism for describing top-down parsers.
Formally, a PEG is a 4-tuple $(V,\Sigma,R,e_S)$, where $V$ is a
finite set of variables, $\Sigma$ is the alphabet, $R$ is the finite
set of rules, and $e_s$ is the start expression. Each rule $r \in R$ is a
pair $(A,e)$, usually written $A \leftarrow e$, where $A \in V$ and $e$
is a parsing expression. We let the meta-variable $a$ denote an
arbitrary alphabet symbol, $A$ a variable and $e$ a parsing expression.
Following common practice, all meta-variables can appear primed or
sub-scripted. The following context-free grammar defines
the syntax of a parsing expression:
\[
\begin{array}{lcl}
e & \to  & \epsilon \,\mid\, a \,\mid\, A\,\mid\,e_1\:e_2\,
\mid\,e_1\,/\,e_2\,\mid\,e^\star\,\mid\,!\,e\\
\end{array}
\]
As usual for PEGs, the unary operators bind more tightly than sequencing, which
in turn binds more tightly than the ordered choice. That is, the not-predicate
$!\,e$ and the Kleene star $e^\star$ have the highest precedence, followed by
sequencing $e_1\,e_2$, and the ordered choice $e_1 / e_2$ has the lowest; thus
$!\,a\,b\,/\,c$ abbreviates $((!\,a)\,b)\,/\,c$. Parentheses may be used to
override these defaults.

The execution of parsing expressions is defined by an inductively defined
big-step judgment that relates a pair of a parsing expression and an input string
to a result. Throughout, the juxtaposition of two strings denotes their
concatenation. The notation $(e,s) \Rightarrow_G (s_p,s_r)$ denotes that parsing
expression $e$ consumes the prefix $s_p$ from the input string $s$, leaving the
suffix $s_r$, so that $s$ equals the concatenation $s_p\,s_r$. We stress that the
right-hand side $(s_p,s_r)$ is a \emph{pair} of the consumed prefix and the
remaining suffix, and not a single string; the rules are big-step, so a
sub-expression occurring in a premise is fully evaluated before the conclusion is
formed. The notation $(e,s) \Rightarrow_G \bot$ denotes the fact that $s$ cannot be
parsed by $e$. We let meta-variable $r$ denote an arbitrary parsing result, i.e.,
either $r$ is a pair $(s_p,s_r)$ or $\bot$. We say that an expression $e$
fails if its execution over an input produces $\bot$; otherwise, it succeeds.
Figure~\ref{fig:pegsemantics} defines the PEG semantics.
\begin{figure*}[h!]
   \[
      \begin{array}{cc}
         \infer[_{\{Eps\}}]{(\epsilon,s)\Rightarrow_G (\epsilon,s)}{} &
         \infer[_{\{ChrS\}}]{(a,as_r) \Rightarrow_G (a,s_r)}{} \\ \\
         \infer[_{\{ChrF\}}]{(a,bs_r) \Rightarrow_G \bot}{a \neq b} &
         \infer[_{\{Var\}}]{(A,s) \Rightarrow_G r}
                           {A \leftarrow e \in R & (e,s) \Rightarrow_G r} \\ \\
         \infer[_{\{Not_S\}}]{(!\,e,s) \Rightarrow_G (\epsilon,s)}
         {(e,s) \Rightarrow_G \bot}
           &
         \infer[_{\{ChrNil\}}]{(a,\epsilon) \Rightarrow_G \bot}{}\\ \\

         \multicolumn{2}{c}{
            \infer[_{\{Cat_{S1}\}}]{(e_1\,e_2,s_{p_1}s_{p_2}s_r) \Rightarrow_G (s_{p_1}s_{p_2}, s_r)}
                                 {(e_1,s_{p_1}s_{p_2}s_r) \Rightarrow_G (s_{p_1},s_{p_2}s_r) &
                                 (e_2,s_{p_2}s_r)\Rightarrow_G (s_{p_2},s_r)}
         } \\ \\
         \multicolumn{2}{c}{
            \infer[_{\{Cat_{F2}\}}]{(e_1\,e_2,s_ps_r) \Rightarrow_G \bot}
                                 { (e_1,s_ps_r) \Rightarrow_G (s_p,s_r) &
                                    (e_2,s_r) \Rightarrow_G \bot}} \\ \\
         \infer[_{\{Cat_{F1}\}}]{(e_1\,e_2,s)\Rightarrow_G \bot}{(e_1,s) \Rightarrow_G \bot} &
         \infer[_{\{Alt_{S1}\}}]{(e_1\,/\,e_2,s_p\,s_r) \Rightarrow_G (s_p,s_r)}
                                {(e_1,s_p\,s_r)\Rightarrow_G (s_p,s_r)}\\ \\
         \multicolumn{2}{c}{
            \infer[_{\{Alt_{S2}\}}]{(e_1\,/\,e_2,s_p\,s_r) \Rightarrow_G r}
                                  {(e_1,s_p\,s_r)\Rightarrow_G \bot &
                                   (e_2,s_p\,s_r)\Rightarrow_G r}
         } \\ \\
         \multicolumn{2}{c}{
            \infer[_{\{Star_{rec}\}}]{(e^\star,s_{p_1}s_{p_2}s_r) \Rightarrow_G (s_{p_1}s_{p_2},s_r)}
                                 {(e,s_{p_1}s_{p_2}s_r) \Rightarrow_G (s_{p_1},s_{p_2}s_r) &
                                  (e^\star, s_{p_2}s_r) \Rightarrow_G (s_{p_2},s_r)}
         }\\ \\
         \multicolumn{2}{c}{
            \infer[_{\{Star_{end}\}}]{(e^\star,s) \Rightarrow_G (\epsilon,s)}
                                    {(e,s) \Rightarrow_G \bot}} \\ \\
         \multicolumn{2}{c}{
            \infer[_{\{Not_F\}}]{(!\,e,s_p\,s_r) \Rightarrow_G \bot}
                               {(e,s_p\,s_r) \Rightarrow_G (s_p,s_r)}
         }
     \end{array}
   \]
   \centering
   \caption{Parsing expressions operational semantics.}
   \label{fig:pegsemantics}
\end{figure*}
We comment on some rules of the semantics. Rule $_{Eps}$ specifies that
expression $\epsilon$ will not fail on any input $s$ by leaving it unchanged.
Rule $_{ChrS}$ specifies that an expression $a$ consumes the first character when
the input string starts with an `a' and rule $_{ChrF}$ shows that it fails when
the input starts with a different symbol. Rule $_{Var}$ parses the input using
the expression associated with the variable in the grammar $G$. When parsing a
sequence expression, $e_1\:e_2$, the result is formed by $e_1$ and $e_2$ parsed
prefixes and the remaining input is given by the result $e_2$. Rules $_{Cat_{F1}}$ and
$_{Cat_{F2}}$ say that if $e_1$ or $e_2$ fail, then the whole expression fails.
The rules for choice impose that we only try
expression $e_2$ in $e_1 / e_2$ when $e_1$ fails. Parsing a star
expression $e^\star$ consists in repeatedly execute $e$ on the input string.
When $e$ fails, $e^\star$ succeeds without consuming any symbol of the input
string. Finally, the rules for the not-predicate expression, $!\,e$, specify
that whenever the expression $e$ succeeds on input $s$, $!\,e$ fails; and when $e$
fails on $s$ we have that $!\,e$ succeeds without consuming any input.

\subsection{A type system for PEGs}

A key issue in the PEG formalism is to ensure the termination of the parsing algorithm
for arbitrary inputs. Ford~\cite{Ford2004} proposed a well-formedness predicate
which is computed as fixed-point over the set of grammar's sub-expressions.
In simple terms, a PEG can loop indefinitely over an input string if: 1) it has
left-recursive rules (whether direct or indirect); 2) expressions that succeed without
consuming any symbol under the Kleene star operator.
As pointed by~\cite{Ribeiro19}, one inconvenience of Ford's well-formedness predicate
is that it is not compositional.
Based on this observation, the work~\cite{Ribeiro19} propose a type
system that guarantees termination by restricting valid expressions to those
that do not obey conditions 1 and 2. In order to ensure these restrictions,
types are interpreted as records containing two fields:
a boolean value indicating that the current expression may succeed without
consuming anything from the input string and a set of variables that can
appear as the first symbol on the recognition branch from the current parsing
expression. This set (named head-set) is defined by recursion on the parsing
expression syntax as follows:
\[
  \begin{array}{lcl}
    head (\epsilon) & = & \emptyset \\
    head (a) & = & \emptyset \\
    head (A) & = & \{A\} \cup head(R(A)) \\
    head (e_1\,e_2) & = & \left\{
                          \begin{array}{ll}
                            head(e_1) \cup head(e_2) & \text{if }e_1 \rightharpoonup 0\\
                            head(e_1) & \text{otherwise}
                          \end{array}
                                        \right. \\
    head(e_1\,/\,e_2) & = & head(e_1) \cup head(e_2) \\
    head(e^\star) & = & head(e) \\
    head(!\,e)    & = & head(e)
  \end{array}
\]
Basically, the head-set of $\epsilon$ and of a terminal $a \in \Sigma$ is empty,
since neither begins with a non terminal. For a non terminal $A$, the head-set is
formed by $A$ itself together with the head-set of its associated parsing
expression $R(A)$; recall that $R$ is the rule function of the grammar, so $R(A)$
is the body of $A$. The head-set of a sequence expression $e_1\,e_2$ includes the
head-set of $e_2$ only when $e_1$ can succeed without consuming any input, a
property we write $e_1 \rightharpoonup 0$ and read as ``$e_1$ is nullable'': in
that case $e_2$ is also reached at the leftmost position. Otherwise
$head(e_1\,e_2)$ is just $head(e_1)$. For the choice operator the head-set is the
union of the head-sets of the two alternatives, since either may be tried first.
Finally, both $e^\star$ and $!\,e$ have the \emph{same} head-set as $e$: each one
begins by attempting to parse $e$ at the current position (the star before deciding whether to iterate, the not-predicate before deciding whether to fail), so the non terminals that may occur at their head are exactly those of $e$. The
two operators differ in whether and how they consume input, not in what they can
start with, which is why $head(!\,e) = head(e^\star) = head(e)$.
In this system the \emph{type} of a parsing expression is precisely the pair of
static informations just computed: whether the expression can succeed without
consuming input, and which non terminals may appear at its head. Following
\cite{Ribeiro19}, we write a type as $\tau = \langle b, S \rangle$, where the
boolean $b$ records nullability (it is $true$ if and only if the expression can succeed consuming no input), and the set $S$ is the head-set. We write
$\tau.null$ for the boolean field $b$ and $\tau.head$ for the head-set $S$. We
call these pairs types because they abstract exactly the aspect of an
expression's behaviour that matters for well-formedness (can it be empty, and what
can it start with), just as a conventional type abstracts the values a program may
produce. The type system manipulates types through the following operations:
\[
  \begin{array}{l}
    b \Rightarrow S = \text{if } b \text{ then } S \text{ else }\emptyset\\
    \tau_1 \otimes \tau_2 = \langle \tau_1.null \land \tau_2.null,
                                \tau_1.head \cup (\tau_1.null \Rightarrow \tau_2.head) \rangle \\
    \tau_1 \oplus \tau_2 = \langle \tau_1.null \lor \tau_2.null,
                                \tau_1.head \cup \tau_2.head\rangle\\
    !\,\tau = \langle true , \tau.head \rangle\\
    \langle false, S \rangle^\star = \langle true, S \rangle
  \end{array}
\]
Rather than restating each equation, we explain why it is defined the way it is.
The guarded operation $b \Rightarrow S$ returns $S$ when $b$ holds and
$\emptyset$ otherwise; its purpose is to make head-set propagation
\emph{conditional} on nullability, and it is the type-level counterpart of the
side condition $e_1 \rightharpoonup 0$ in the definition of $head$. The product
$\tau_1 \otimes \tau_2$ is the type of a sequence $e_1\,e_2$: the sequence is
nullable only if both parts are (hence the conjunction), and its head-set is that
of $e_1$ extended with that of $e_2$ \emph{exactly when} $e_1$ is nullable and so
$e_2$ can be reached first, which is why $\tau_2.head$ is guarded by
$\tau_1.null$. The co-product $\tau_1 \oplus \tau_2$ is the type of a choice
$e_1 / e_2$: either alternative may match, so the result is nullable if
\emph{either} operand is (a disjunction), and its head-set is the union of both.
The not operation $!\,\tau$ marks the type nullable, because a not-predicate never
consumes input and so always succeeds on the empty string, while keeping the
head-set since $!\,e$ still inspects $e$ at the head. The last operation is the
Kleene star on types: here we overload the notation $(\cdot)^\star$, defined
earlier on expressions, to act on type pairs. It also sets the boolean field to
$true$, since a star can always match zero repetitions; crucially, it is defined
\emph{only} for non-nullable types ($\tau.null = false$) and left undefined
otherwise. This is deliberate: a star over a nullable body could iterate forever
without consuming input, so leaving $(\cdot)^\star$ undefined on nullable types is
exactly how the type system forbids that source of non-termination. Finally, following the common practice, the meta-variable
$\Gamma$ denotes a typing context and notation $\Gamma(A) = \tau$ holds whenever
$(A,\tau) \in \Gamma$.
The type system is defined as an inductive judgment
$\Gamma \vdash e : \tau$ in Figure~\ref{fig:typesystem}. A
PEG $G = (V,\Sigma,R,e_S)$ is considered well-typed, written $\Gamma \vdash G$,
if $\forall A \in V . R(A) : \Gamma (A)$.

\begin{figure}[h!]
\[
  \begin{array}{cc}
    \infer{\Gamma \vdash \epsilon : \langle true , \emptyset \rangle}{} &
    \infer{\Gamma \vdash a : \langle false , \emptyset\rangle}{} \\ \\
    \multicolumn{2}{c}{
      \infer{\Gamma \vdash A : \langle \tau.null, \tau.head \cup {A}\rangle}
            {\Gamma(A) = \tau & A \not\in \tau.head}
    }\\ \\
    \infer{\Gamma\vdash e_1\,/\,e_2 : \tau_1 \oplus \tau_2}
          {\Gamma \vdash e_1 : \tau_1 & \Gamma \vdash e_2 : \tau_2}
    &
    \infer{\Gamma\vdash !\,e :\, !\,\tau}
          {\Gamma \vdash e : \tau}\\ \\
    \infer{\Gamma\vdash e_1\,e_2 : \tau_1 \otimes \tau_2}
          {\Gamma \vdash e_1 : \tau_1 & \Gamma \vdash e_2 : \tau_2} &
    \infer{\Gamma\vdash e^\star : \tau^\star}
          {\begin{array}{c}
             \Gamma \vdash e : \tau\\
             \tau.null = false\\
           \end{array}}
  \end{array}
\]
\centering
\caption{Type system definition.}
\label{fig:typesystem}
\end{figure}

The first rule of the type system specifies that a parsing
expression $\epsilon$ has type $\langle true, \emptyset \rangle$.
The rule of a single character expression says that its type is
$\langle false, \emptyset\rangle$, since it does not succeed without
consuming a symbol and no variable can appear as the head of its
recognition branch, i.e., its head-set is $\emptyset$. The type of a prioritized
choice expression ($e_1/e_2$) has type $\tau_1 \oplus \tau_2$,
where $\Gamma \vdash e_1 : \tau_1$ and $\Gamma \vdash e_2 : \tau_2$.
The type of a sequence expression is given by the product operation
on their sub-expressions types, and the star operation is only well-typed
if its underlying expression does not succeed without consuming
a symbol. A not-predicate parsing expression always has the type
$\langle true,\tau.head \rangle$ whenever its sub-expression is well-typed.
Finally, a variable is well-typed only if its type is in the context and
itself does not belong to its head-set.
The type system proposed by~\cite{Ribeiro19} readily rejects invalid expressions like
$a(!b / c)^\star$, since it only considers well-typed a Kleene parsing
expressions in which an underlying expression does
not succeed without consuming the input. However, an inconvenience of it is
the need to require type annotations on all grammar's non terminals.
A later work~\cite{Cardoso23} proposes a type inference algorithm, which
guarantees termination of the parsing process.

\subsection{An overview of the Agda programming language}
We mechanise our metatheory in Agda~\cite{agda}, a dependently-typed functional
language and proof assistant based on Martin-L\"of type theory. Agda's totality
discipline, namely that every function must be shown to cover all inputs and to
terminate, is what makes it a natural setting for our result: a proof of strong
normalization is, quite literally, a total function that produces a parse result
for every well-typed expression on every input. This subsection introduces,
through small examples, the handful of features we rely on.

A data type is introduced by listing its constructors, and a function over it is
defined by pattern matching and structural recursion. The natural numbers and
their addition are the prototypical example.
\begin{code}[hide]
module AgdaNat where
  open import Relation.Binary.PropositionalEquality using (_≡_; refl; cong)
\end{code}
\clearpage
\begin{code}
  data ℕ : Set where
    zero : ℕ
    suc  : ℕ → ℕ

  _+_ : ℕ → ℕ → ℕ
  zero + n = n
  suc m + n = suc (m + n)
\end{code}
In a dependently-typed setting propositions are types, and proofs are programs
that inhabit them. Propositional equality $x \equiv y$ is itself a data type whose
only constructor \lstinline{refl} proves $x \equiv x$. A proof by induction is
then a recursive function: to show that \lstinline{zero} is a right identity of
addition we pattern match on the argument, returning \lstinline{refl} in the base
case and appealing to the induction hypothesis in the step case, where
\lstinline{cong} applies \lstinline{suc} congruently to it.
\begin{code}
  +-identityʳ : ∀ (n : ℕ) → n + zero ≡ n
  +-identityʳ zero = refl
  +-identityʳ (suc n) = cong suc (+-identityʳ n)
\end{code}

A data type may also be \emph{indexed} by values, which is the essence of
dependent types. The finite-set type \lstinline{Fin n}, which we use pervasively
in the formalization to name the non-terminals of a grammar, is the standard
example. It is indexed by a natural number and represents the finite set
$\{0, 1, \dots, n-1\}$.
\begin{code}[hide]
module AgdaFin where
  open import Data.Nat using (ℕ; suc)
\end{code}
\begin{code}
  data Fin : ℕ → Set where
    zero : ∀ {n} → Fin (suc n)
    suc  : ∀ {n} → Fin n → Fin (suc n)
\end{code}
The index makes the constructors precise: \lstinline{zero} inhabits every
non-empty \lstinline{Fin (suc n)}, and \lstinline{suc} embeds \lstinline{Fin n}
into \lstinline{Fin (suc n)}. In particular \lstinline{Fin 0} is empty, since
neither constructor can build it, and \lstinline{Fin n} has exactly $n$ elements.
In the formalization a grammar with $n$ non-terminals uses \lstinline{Fin n} to
index its rule table, and a head set is a \emph{subset} of \lstinline{Fin n}; the
cardinality of such a subset is precisely the measure that the well-founded
recursion below decreases.

Head sets are values of the library type \lstinline{Subset n}, a vector of $n$ bits
recording which elements of \lstinline{Fin n} the set contains. From that module we use the empty set
$\bot$, the singleton $\lceil i \rceil$, the union $p \cup q$, membership
$x \in p$ and its negation (acyclicity of a grammar is the statement that a
non-terminal does not belong to its own head set), the inclusion
$p \subseteq q$, its strict version $p \subset q$, which is an inclusion
together with a witness lying in $q$ but not in $p$, and the cardinality $|p|$.
Because the representation is a bit vector, membership and cardinality are
computable, and it is the cardinality that our recursion decreases.

Besides routine membership and union lemmas, \lstinline{Data.Fin.Subset.Properties}
supplies the two results that carry the termination argument. We display them
under descriptive names; the definitions are literal aliases, so the types shown
are the ones the library proves.
\begin{code}[hide]
module AgdaSubset where
  open import Data.Nat using (ℕ; _<_; _≤_)
  open import Data.Fin using (Fin)
  open import Data.Fin.Subset using (Subset; _⊆_; _⊂_; ∣_∣)
  open import Data.Fin.Subset.Properties using (p⊂q⇒∣p∣<∣q∣; p⊆q⇒∣p∣≤∣q∣)
\end{code}
\begin{code}
  size-mono   : ∀ {n} {p q : Subset n} → p ⊆ q → ∣ p ∣ ≤ ∣ q ∣
  size-mono   = p⊆q⇒∣p∣≤∣q∣

  size-strict : ∀ {n} {p q : Subset n} → p ⊂ q → ∣ p ∣ < ∣ q ∣
  size-strict = p⊂q⇒∣p∣<∣q∣
\end{code}
The first propagates a measure bound from a compound expression down to its
sub-expressions. The second is what makes the recursion terminate: expanding a
non-terminal $i$ replaces its head set $\lceil i \rceil \cup S$ by $S$, where
$S$ is the head set assigned to $i$ by the context, and acyclicity supplies the
witness $i$, which lies in the former but not in the latter. The inclusion is
therefore strict, and \lstinline{size-strict} turns it into a strict decrease of
a natural number, on which the well-founded recursion described next is
performed.

Not every terminating function recurses on a structurally smaller argument. When
the decreasing quantity is a \emph{measure} rather than a subterm, the termination
checker is satisfied by \emph{well-founded recursion}, which is the central
technique of our proof. The accessibility predicate \lstinline{Acc _<_ x} holds
when every \lstinline{y} with $y < x$ is itself accessible; a relation is
well-founded when every element is accessible, and the standard library proves
that the strict order on the naturals, \lstinline{<-wellFounded}, is. One recurses
by pattern matching on the accessibility proof, which supplies, for each smaller
argument, its own accessibility proof. Integer division illustrates the pattern:
dividing \lstinline{n} by \lstinline{suc m} recurses on the truncated difference
of \lstinline{n} and \lstinline{suc m}, which is strictly smaller whenever
$n \ge \mathtt{suc}\ m$.
\begin{code}[hide]
module AgdaDiv where
  open import Data.Nat using (ℕ; zero; suc; _∸_; _<_; _<?_)
  open import Data.Nat.Properties using (∸-monoʳ-<; ≮⇒≥; 0<1+n)
  open import Data.Nat.Induction using (Acc; acc; <-wellFounded)
  open import Relation.Nullary using (yes; no)
\end{code}
\clearpage
\begin{code}
  divsuc : ℕ → ℕ → ℕ           -- divsuc n m is the quotient of n by (suc m)
  divsuc n m = go n (<-wellFounded n)
    where
    go : (k : ℕ) → Acc _<_ k → ℕ
    go k (acc rs) with k <? suc m
    ... | yes _ = zero
    ... | no k≮ = suc (go (k ∸ suc m) (rs (∸-monoʳ-< 0<1+n (≮⇒≥ k≮))))
\end{code}
This is precisely the technique underlying our termination proof. Parsing a
Kleene star recurses on the strictly shorter remaining input, and expanding a
non-terminal recurses on a strictly smaller head set; neither is a structural
subterm, so both are justified by well-founded recursion, and the two measures are
combined lexicographically (Section~\ref{sec:typesoundnesslean}).

\subsection{Strong normalization via logical relations}
Our termination proof follows the classical method that Tait introduced to prove
strong normalization of the simply-typed lambda calculus, and which is by now the
standard tool for such results~\cite{plfa22.08}. We recall its structure here,
since our development is a direct adaptation of it.

The simply-typed lambda calculus has types built from a base type $o$ by the
function-space constructor, and terms that are variables, abstractions and
applications,
\[
  \sigma, \tau \;::=\; o \;\mid\; \sigma \to \tau,
  \qquad\qquad
  t, u \;::=\; x \;\mid\; \lambda x{:}\sigma.\,t \;\mid\; t\,u,
\]
typed by the usual rules under a context $\Gamma$ that assigns types to variables.
Computation is $\beta$-reduction, $(\lambda x{:}\sigma.\,t)\,u \to t[u/x]$, closed
under the term constructors. A term is \emph{strongly normalizing}, written
$\mathrm{SN}(t)$, when every reduction sequence starting from it is finite. The
goal is to show that every well-typed term is strongly normalizing.

A direct induction on the typing derivation fails at the application rule: from
$\mathrm{SN}(t)$ and $\mathrm{SN}(u)$ one cannot conclude $\mathrm{SN}(t\,u)$,
because reducing $t$ to an abstraction and substituting $u$ may create new redexes
that neither premise controls. Tait's insight is to prove something stronger, by
induction on \emph{types} rather than on terms. One defines a type-indexed family
of predicates on terms, the \emph{logical relation} (also called the
\emph{reducibility} predicate) $\mathrm{Red}_\sigma$, by
\[
  \begin{array}{lcl}
    \mathrm{Red}_{o}(t)             & \iff & \mathrm{SN}(t),\\
    \mathrm{Red}_{\sigma\to\tau}(t) & \iff & \text{for all } u,\ \mathrm{Red}_\sigma(u)\ \text{implies}\ \mathrm{Red}_\tau(t\,u).
  \end{array}
\]
At the base type reducibility is just strong normalization; at a function type it
is the requirement that the term send reducible arguments to reducible results.

The proof rests on three properties of this family, the so-called reducibility
candidate conditions.

\begin{Lemma}[Reducibility candidates]\label{lem:cr}
For every type $\sigma$, the predicate $\mathrm{Red}_\sigma$ satisfies:
\begin{enumerate}
  \item[(CR1)] reducibility implies strong normalization, that is,
    $\mathrm{Red}_\sigma(t)$ implies $\mathrm{SN}(t)$;
  \item[(CR2)] reducibility is preserved by reduction, that is, if
    $\mathrm{Red}_\sigma(t)$ and $t \to t'$ then $\mathrm{Red}_\sigma(t')$;
  \item[(CR3)] a \emph{neutral} term (a variable or an application, that is, one
    that is not an abstraction) all of whose one-step reducts are reducible is
    itself reducible. In particular, every variable is reducible.
\end{enumerate}
\end{Lemma}

\begin{proofsketch}
The three properties are proved simultaneously by induction
on the type. At the base type they follow from the definition and from basic
facts about $\mathrm{SN}$. At a function type, (CR1) applies the term to a
reducible variable, available by (CR3) at $\sigma$, and then uses (CR1) at
$\tau$; (CR2) is immediate; and (CR3) uses (CR2) at $\sigma$ and (CR1) with
(CR3) at $\tau$ to account for the redexes an application may create.
\end{proofsketch}

The heart of the argument is the fundamental theorem, which states that every
well-typed term is reducible under any reducible substitution.

\begin{Theorem}[Fundamental theorem]\label{thm:fundamental-stlc}
If $\Gamma \vdash t : \sigma$ and $\rho$ maps each variable $x{:}\tau$ of
$\Gamma$ to a term with $\mathrm{Red}_\tau(\rho(x))$, then
$\mathrm{Red}_\sigma(t[\rho])$.
\end{Theorem}

\begin{proofsketch}
By induction on the typing derivation. The variable case is
the hypothesis on $\rho$; the application case is now immediate, being exactly
the definition of reducibility at a function type; and the abstraction case uses
(CR3) together with (CR1) and (CR2) to show that a lambda sends reducible
arguments to reducible results.
\end{proofsketch}

\begin{Theorem}[Strong normalization]\label{thm:sn-stlc}
If $\Gamma \vdash t : \sigma$ then $\mathrm{SN}(t)$.
\end{Theorem}

\begin{proofsketch}
Instantiate Theorem~\ref{thm:fundamental-stlc} with the
identity substitution, which is reducible because variables are reducible by
(CR3). This gives $\mathrm{Red}_\sigma(t)$, and hence $\mathrm{SN}(t)$ by (CR1).
\end{proofsketch}

Our development in Section~\ref{sec:typesoundnesslean} follows this template. The
logical relation becomes the predicate $\mathcal{R}$, which pairs
\emph{termination} of parsing (the analogue of strong normalization, since a total
big-step relation plays the role of a terminating reduction) with a
\emph{progress} property (the analogue of a reduction step making progress through
the input); the closure lemmas play the role of (CR1) to (CR3); and the
fundamental lemma is the direct counterpart of the fundamental theorem. The one
ingredient with no analogue in the simply-typed lambda calculus is recursion
through the grammar: a non-terminal expands to its rule body, which is not a
subterm, so the induction that is purely structural in the lambda calculus must
here be strengthened by the well-founded induction on the head-set rank.

\section{Formalizing Pest operators}~\label{sec:pest-formalization}

% We start by presenting the syntax of the new PEG operators, together with
% some examples. Next, we present the formal semantics of the new operators,
% giving details on repetition indexes are normalized. Finally, we describe
% our extension of Ribeiro et. al. type system~\cite{Ribeiro19} which
% guarantee termination of grammars using the new operators.

\paragraph{New operators syntax}
The Pest tool includes two types of operators: stack and repetition operators.
The syntax of these operators are defined next:
\[
\begin{array}{lcl}
  e & \to  & ...\,\mid\,e^i\,\mid\,e[i,i]\,\mid\,\mathbf{push}\,\mid\,\mathbf{pop}\,\mid\,\mathbf{peek}
             \,\mid\,\mathbf{drop}\\
    & \mid & \mathbf{peekall}\,\mid\,\mathbf{dropall}\\
  i & \to  & \mathbf{top.tonat}\,\mid\,\mathbf{top.length}\,\mid\,i \circ i\,\mid\,\nu\\
  \nu & \to  & \mathbf{infty}\,\mid\,n,m\in\mathbb{N}\\
  \sigma & \to & s : \sigma\,\mid\,[] \\
  \circ & \to & + \,\mid\, \times
\end{array}
\]
Meta-variable $s$ denotes an arbitrary string using grammar's alphabet.
We let meta-variable $\sigma$ denote the parser stack and we represent it
using Haskell-style list notation.

\paragraph{Semantics of new operators}
Now that we have introduced the syntax and the informal semantics of the new Pest
operators, we turn our attention to its formal semantics. One important aspect
of Pest repetition operators is the usage of indexes. Before evaluating any
of these new operators we need to reduce indexes to a normal form.
We consider that an index is in normal form if it is a number or the constant
\textbf{infty}. We let meta-variable $v$ denote an arbitrary index normal form.
In simple terms, the index normalization algorithm performs the following steps:

\begin{enumerate}
  \item Evaluate $\mathbf{top.length}$ and $\mathbf{top.tonat}$, if
    they are present in the current index. The usage of these operators
    can result in a run-time error if the parser stack is empty or the
    stack first element can not be converted to a number (for $\mathbf{top.tonat}$).
  \item Reduce all arithmetic operations using their usual meaning and
    the following reduction rules
    when the constant \textbf{infty} is present:
\[
  \begin{array}{c}
    n \circ \mathbf{infty} \rightarrow \mathbf{infty}\\
    \mathbf{infty}\circ n \rightarrow \mathbf{infty}\\
    \mathbf{infty}\circ\mathbf{infty} \rightarrow \mathbf{infty}\\
  \end{array}
\]
where $\circ$ denotes either addition or multiplication and $n$ a
numeric value.
\end{enumerate}

The semantics is denoted by the inductive judgment
(Figures~\ref{fig:newsemantics},~\ref{fig:exactsemantics} and~\ref{fig:intervalsemantics})
$(e, w, \sigma) \Rightarrow_G r$,
where $e$ is the current expression being executed, $w$ is the current input string,
$\sigma$ is the current input parser stack, $G$ is the grammar and $r$ is the output
which can be: 1) success, denoted by triple $(w_1,w_2,\sigma')$ where $w_1$
is the parsed prefix, $w_2$ is the remaining input and $\sigma'$ is the resulting
stack from executing the expression $e$, or; 2) failure, denoted by $\bot$.
The inclusion of a stack requires changing the semantics of all standard PEG
operators to thread the current stack state through parser execution.
We omit these modified rules, since they don't add any interesting insight.
\begin{figure*}[h!]
  \[
    \begin{array}{cc}
      \infer[_{\{PUSH_S\}}]
             {(\mathbf{push}\:e,w,\sigma) \Rightarrow_G (w_1, w_2, (w_1:\sigma'))}
             {(e,w,\sigma)\Rightarrow_G (w_1,w_2,\sigma')} &
      \infer[_{\{POP_S\}}]
            {(\mathbf{pop},w,w_1\,:\,\sigma)\Rightarrow_G (w_1,w_2,\sigma)}
            {\exists w_2. w = w_1w_2} \\ \\
        \infer[_{\{DROP_F\}}]
              {(\mathbf{drop},w,[])\Rightarrow_G \bot}{} &
      \infer[_{\{POP_F\}}]
            {(\mathbf{pop},w,\sigma)\Rightarrow_G \bot}
            {\neg \exists w_1\,w_2\, \sigma'. w = w_1w_2 \land \sigma = w_1\,:\,\sigma'}\\ \\
      \multicolumn{2}{c}{
        \infer[_{\{PEEK_F\}}]
              {(\mathbf{peek},w,\sigma)\Rightarrow_G \bot}
              {\neg \exists w_1\,w_2\, \sigma'. w = w_1w_2 \land \sigma = w_1\,:\,\sigma'}
      }\\\\
      \multicolumn{2}{c}{
        \infer[_{\{PEEKALL_S\}}]
              {(\mathbf{peekall},w,\sigma)\Rightarrow_G (w', w'',\sigma)}
              {\sigma = [w_0, ..., w_n] & w' = w_0...w_n & \exists w'' . w = w'w''}
      }\\ \\
      \infer[_{\{PEEK_S\}}]
              {(\mathbf{peek},w,w_1\,:\,\sigma)\Rightarrow_G (w_1,w_2,w_1 \,:\,\sigma)}
              {\exists w_2. w = w_1w_2} &
      \infer[_{\{DROP_S\}}]
            {(\mathbf{drop},w,w_1\,:\,\sigma)\Rightarrow_G (\epsilon,w,\sigma)}
            {}\\ \\
      \multicolumn{2}{c}{
        \infer[_{\{PEEKALL_F\}}]
              {(\mathbf{peekall},w,\sigma)\Rightarrow_G \bot}
              {\sigma = [w_1, ..., w_n] & w' = w_1...w_n & \neg \exists w'' . w = w'w''}
      }\\ \\
         \infer[_{\{PUSH_F\}}]
              {(\mathbf{push}(e),w,\sigma)\Rightarrow_G \bot}
              {(e,w,\sigma)\Rightarrow_G \bot} &
             \infer[_{\{DROPALL_F\}}]
                   {(\mathbf{dropall}, w, [])\Rightarrow_G \bot}
                   {}
          \\ \\
          \multicolumn{2}{c}{
           \infer[_{\{DROPALL_S\}}]
                 {(\mathbf{dropall}, w, w_1 : \sigma)\Rightarrow_G (\epsilon, w, [])}
                 {}
          }
  \end{array}
  \]
  \caption{Semantics of the stack operators}
  \label{fig:newsemantics}
\end{figure*}

Figure~\ref{fig:newsemantics} presents the semantics for the stack manipulation operators.
Rules $PUSH_S$ and $PUSH_F$ formalize the behaviour of the push operator, which pushes the
prefix parsed by the push's expression argument into the parser stack or it fails when
the expression fails. The semantics of the \textbf{pop} operator is defined by rules
$POP_S$ and $POP_F$, which removes the topmost element of the stack and matching it
against the input string. When such a matching is not possible, we have a failure.
The only difference between the semantics for the \textbf{peek} operation and \textbf{pop}
is that the former keeps the top element in the result stack.
The \textbf{drop} operator removes the top element of the stack without matching it
against the current input or it fails when the stack is empty (check rules $DROP_S$ and
$DROP_F$). Rules $PEEKALL_S$ and $PEEKALL_F$ use the whole stack content to match
the current input, failing if all the stack elements do not form a prefix of the current
input and $DROPALL_S$ simply empties the parser stack.

\begin{figure*}[h!]
  \[
    \begin{array}{c}
      \infer[_{\{EX_{S1}\}}]
            {(e^i, w, \sigma)\Rightarrow_G (\epsilon,w,\sigma)}
            {\mathbf{norm}(i,\sigma) = 0} \\ \\
        \infer[_{\{EX_{S2}\}}]
              {(e^i,w,\sigma) \Rightarrow_G (w_1w'_1,w'_2,\sigma')}
              {\begin{array}{c}
                \mathbf{norm}(i,\sigma) = n \\
                (e,w,\sigma)\Rightarrow_G(w_1,w_2,\sigma_1)
               \end{array} &
               \begin{array}{c}
                 0 < n \\
                 (e^{n - 1},w_2,\sigma_1)\Rightarrow_G (w'_1,w'_2,\sigma')
               \end{array}}\\ \\
          \infer[_{\{EX_{F1}\}}]
                {(e^i,w,\sigma)\Rightarrow_G \bot}
                {\begin{array}{c}
                  \mathbf{norm}(i,\sigma) = n \\
                  0 < n \quad (e,w,\sigma)\Rightarrow_G \bot
                \end{array}} \\ \\
            \infer[_{\{EX_{F2}\}}]
                  {(e^i,w,\sigma) \Rightarrow_G \bot}
                  {\begin{array}{c}
                    \mathbf{norm}(i,\sigma) = n \\
                    (e,w,\sigma)\Rightarrow_G(w_1,w_2,\sigma_1)
                  \end{array} &
                  \begin{array}{c}
                    0 < n \\
                    (e^{n - 1},w_2,\sigma_1)\Rightarrow_G \bot
                  \end{array}}
    \end{array}
  \]
  \caption{Semantics for the exact repetition operator}
  \label{fig:exactsemantics}
\end{figure*}

The semantics for the exact repetition operator is presented in
Figure~\ref{fig:exactsemantics}.
The first rule, $EX_{S1}$, shows that for any expression $e$, $e^0$ consumes the
empty string from the input and does not change the parser stack. The second rule,
$EX_{S2}$, shows that to evaluate $e^i$, we first normalize the index $i$ to obtain
a numeric value, $n$. Next, we check if $0 < n$ and parse the input using $e$ producing
a remaining suffix $w_2$, which is used as input for running the expression $e^{n - 1}$.
We say that the repetition operator $e^i$ fails if any iteration of $e$ results in a failure,
as formalized by rules $EX_{F1}$ and $EX_{F2}$.


\begin{figure*}[h!]
  \[
    \begin{array}{cc}
      \infer[_{\{Inter_{1}\}}]
            {(e[i,j],w,\sigma) \Rightarrow_G r}
            {\begin{array}{c}
               \mathbf{norm}(i,\sigma) = n \\
               \mathbf{norm}(j,\sigma) = n \\
               (e^n,w, \sigma)\Rightarrow r
           \end{array}} &
      \infer[_{\{Inter_2\}}]
            {(e[i,j],w,\sigma)\Rightarrow_G (w_1,w_2,\sigma')}
            {\begin{array}{c}
              \mathbf{norm}(i,\sigma) = n \\
              \mathbf{norm}(j,\sigma) = m \\
              n < m \\
              (e^{m}, w, \sigma)\Rightarrow_G (w_1,w_2,\sigma')
          \end{array}} \\ \\
      \infer[_{\{Inter_{F1}\}}]
            {(e[i,j], w, \sigma) \Rightarrow_G \bot}
            {\begin{array}{c}
                \mathbf{norm}(i,\sigma) = n\\
                \mathbf{norm}(j,\sigma) = m\\
                n > m\\
            \end{array}} &
        \infer[_{\{Inter_3\}}]
              {(e[i,j], w, \sigma) \Rightarrow_G r}
              {\begin{array}{c}
                  \mathbf{norm}(i,\sigma) = n \\
                  \mathbf{norm}(j,\sigma) = m \\
                  n < m \\
                  (e^{m},w,\sigma)\Rightarrow_G \bot  \\
                  (e[n,m - 1],w, \sigma)\Rightarrow_G r
              \end{array}}
        \\ \\
        \multicolumn{2}{c}{
          \infer[_{\{Inter_{I}\}}]
                {(e[i,j], w, \sigma) \Rightarrow_G r}
                {\begin{array}{c}
                    \mathbf{norm}(i,\sigma) = n\\
                    \mathbf{norm}(j,\sigma) = \mathbf{infty} \\
                    (e^n\,e^{\star},w,\sigma)\Rightarrow_G r
                  \end{array}
              }}
    \end{array}
  \]
  \centering
  \caption{Semantics for the interval repetition operator}
  \label{fig:intervalsemantics}
\end{figure*}

The last set of semantics rules deals with interval repetition operator,
$e[i,j]$. Rule $Inter_{1}$ parses the input using $e^n$, when both of indexes
normalize to the same numeric value $n$. The rule $Inter_2$ shows how to run
an interval repetition when the normalized indexes are such that $n < m$,
where $\mathbf{norm}(i,\sigma) = n$ and $\mathbf{norm}(j,\sigma) = m$:
it tries to parse the input using $e^m$, if it succeeds then the result is
returned. In case of failure of $e^m$, rule $Inter_3$ tries
to parse the input using the expression $e[n, m -1]$.
Rule $Inter_{F1}$ specifies that an interval repetition fails when the formalized value of $i$
is greater than normalized value of $j$.
The rule deals with the situation when the upper limit bound normalizes to \textbf{infty}.
In this situation, expression $e[i,j]$ is interpreted as $e^{n}\,e^{\star}$, where $n$ is the
result of normalizing the index $i$.

\subsection{Type system and its inference algorithm}~\label{sec:typesystem}

Now, we turn our attention to extending a type system to avoid
non termination problems during a grammar execution. We follow
the approach of defining typing rules and its corresponding
type inference using the strategy of constraint generation
followed by a solving step~\cite{Ribeiro19,Cardoso23, Pottier05}.

Before starting discussing the typing rules for the new operators,
we need to extend the $head$ function. In Figure~\ref{fig:new-head-set},
we present the new equations which handle both repetions and stack
operators.

\begin{figure}[H]
  \[
    \begin{array}{lcl}
      head(e^i) & = & head(e)\\
      head(e[i,j]) & = & head(e)\\
      head(\mathbf{push}(e)) & = & head(e)\\
      head(\mathbf{pop}) & = & \emptyset\\
      head(\mathbf{peek}) & = & \emptyset\\
      head(\mathbf{drop}) & = & \emptyset\\
      head(\mathbf{dropall}) & = & \emptyset\\
      head(\mathbf{peekall}) & = & \emptyset\\
    \end{array}
  \]
  \caption{Head-set equations for the new operators}
  \label{fig:new-head-set}
\end{figure}
The first three equations shows that the head-set of
repetitions and the push operator are equal to their
parameter expression set. All the remaining stack operators
have their head-set equal to $\emptyset$, since they do not
refer to non terminals.

The following subsections present the extensions needed to
support typing the new operators.

\subsubsection{Typing rules for the new operators}

We start by defining the typing rules for our new PEG repetition operators.
Figure~\ref{fig:typing-repetition} presents the typing rules for the new
repetition operators. Typing rules make use of the following notations:
1) $\mathbf{norm}^{\ast}(i)$ denotes a static analysis
version of the normalization algorithm in which operators \textbf{top.tonat}
and \textbf{top.length} are replaced by 0\footnote{This is necessary because the value
of these operators depends on the current parser stack state, which is unknown
during type inference.}; 2) we let $\tau^\nu$ denote the
type $\langle \nu = 0 \lor \tau.null, \tau.head\rangle$ and
3) notation $\tau[\nu_1,\nu_2]$ denote the type
$\langle \nu_1 = 0 \lor \tau.null, \tau.head\rangle$.


The first type system rule shows that an exact repetition is considered
nullable if either its underlying expression is nullable or its normalized index
is equal to $0$. The head-set for an exact repetition is
the same of its parameter expression. We also restrict that normalized indexes can
only be numeric values, $n$.
Typing of the interval repetition operator is similar to exact repetition, but
we allow the upper interval bound to be \textbf{infty}.
Intervals are considered nullable when
its parameter expression is nullable or either the normalized left index, $n$,
is equal to zero. We also do not allow intervals that have an \textbf{infty}
upper bound to be nullable, since such parsing expression will diverge
whenever they succeed without consuming any input.

\begin{figure}[H]
  \[
    \begin{array}{cc}
      \infer{\Gamma \vdash e^i : \tau^n}
            {\begin{array}{c}
              \mathbf{norm}^{\ast}(i) = n\quad \Gamma \vdash e : \tau \\
          \end{array}} &
        \infer{\Gamma\vdash e[i,j] : \tau[n,\nu] }
              {\begin{array}{ccc}
                  \Gamma \vdash e : \tau &
                  \mathbf{norm}^{\ast}(i) = n &
                  \mathbf{norm}^{\ast}(j) = \nu \\
                  \multicolumn{3}{c}{
                  \neg (\nu \equiv \mathbf{infty} \land\tau.null)}
               \end{array}}
    \end{array}
  \]
  \caption{Typing rules for new PEG repetition operators}
  \label{fig:typing-repetition}
\end{figure}

Figure~\ref{fig:typing-stack} presents the typing rules for the PEG stack
operators. The first rule specifies that the type for \textbf{push}($e$) is
the same as $e$, since it simply pushes the $e$'s consumed input and its
head-set will be the same as $e's$ head-set. Since all other stack operators
can succeed without consuming input, all of them are considered nullable.
As presented in Figure~\ref{fig:new-head-set}, the head-set of \textbf{pop},
\textbf{drop}, \textbf{peek}, \textbf{peekall} and \textbf{dropall} are the
empty set.

\begin{figure}[H]
  \[
    \begin{array}{ccc}
      \infer{\Gamma \vdash \mathbf{push}(e) : \tau}
            {\Gamma \vdash e : \tau}
      &
      \infer{\Gamma \vdash \mathbf{pop} : \langle true, \emptyset\rangle}{} &
      \infer{\Gamma \vdash \mathbf{peek} : \langle true, \emptyset\rangle}{} \\ & \\
      \infer{\Gamma \vdash \mathbf{peekall} : \langle true, \emptyset\rangle}{} &
      \infer{\Gamma \vdash \mathbf{drop} : \langle true, \emptyset\rangle}{} &
      \infer{\Gamma \vdash \mathbf{dropall} : \langle true, \emptyset\rangle}{}
    \end{array}
  \]
  \caption{Typing rules for new PEG stack operators}
  \label{fig:typing-stack}
\end{figure}

\subsection{Extending the type inference algorithm}

In this section we describe the necessary extensions to the type inference algorithm
created by Cardoso et. al.~\cite{Cardoso23} to support the new repetitions and
stack operators.

\paragraph{Extending the constraint syntax} We start by including index normal forms
on the constraint grammar (taken from~\cite{Cardoso23}). The non terminal $\nu$ denotes
a normal form of an index expression, which can be a numeric value or
constant \textbf{infty}. The complete constraint grammar is as follows:

  \[
    \begin{array}{lcl}
      \tau & \to & \alpha\,\mid\,\langle b, S \rangle\\
      \nu & \to & n\,\mid\,\mathbf{infty}\\
      \mathbf{t} & \to & A\,\mid\,b\,\mid\,S\,\mid\,\tau\,\mid\,\nu\\
      C    & \to & \mathbf{true}\,\mid\,\mathbf{false}\,\mid\,\mathbf{t} \equiv \mathbf{t}\,\mid\,C \land C\,\mid\, \exists
                   \alpha.C\,\mid\, \mathbf{def}\:A :\tau\:\mathbf{in}\:C
    \end{array}
  \]
Constraints are first-order logic formulas over types and their meaning
can be defined inductively. We omit the constraint semantics, since its
definition is the same presented in~\cite{Cardoso23}.

\paragraph{Extending constraint generation and solving} After including index expressions
in the constraint syntax, we need to include equations for the new operations
in the constraint generation algorithm.
In order to generate constraints for the exact repetition operator, we use
an existential quantifier to define a new variable, $\alpha$, that is passed
as argument to the generation of constraints for $e$. The algorithm creates
an equality between the argument type, $\tau$, and the repetition operator
type over $\alpha$. We also include an inequality to assert that the
normalized index should not be \textbf{infty}.
\[
  \langle\langle e^i : \tau \rangle\rangle = \exists \alpha . \langle\langle e,\alpha\rangle\rangle \land \tau \equiv \alpha^{\mathbf{norm}^{\ast}(i)} \land \mathbf{norm}^{\ast}(i) \not\equiv \mathbf{infty}\\
\]
Next, we consider the generation of constraints for the interval repetition operator, which
follows the same pattern of the exact repetition: we generate the constraint for sub-expression
$e$, considering an existential quantified type variable and create an equality between the variable
type and the input argument type, $\tau$. We also impose that the interval lower-bound
cannot be \textbf{infty} and that if the upper-bound is \textbf{infty}, the interval
underlying expression cannot be nullable since it could make the parser loop indefinitely.
\[
  \begin{array}{rl}
    \langle\langle e[i,j] : \tau \rangle\rangle =   & \exists \alpha . \langle\langle e, \alpha\rangle\rangle \land \tau \equiv \alpha[\mathbf{norm}^{\ast}(i),\mathbf{norm}^{\ast}(j)]\\
                                            \land{} & \mathbf{norm}^{\ast}(i) \not\equiv\mathbf{infty} \\
                                            \land{} & \neg (\mathbf{norm}^{\ast}(j) \equiv \mathbf{infty} \land \alpha.null)\\
  \end{array}
\]
The constraint for the \textbf{push} operator is just the constraint of its argument expression and
all other stack expressions just generate an equality between its input argument type and
$\langle true,\emptyset\rangle$.

\begin{figure}[H]
  \[
    \begin{array}{lcl}
      \langle\langle \mathbf{push}(e) : \tau\rangle\rangle & = & \langle\langle e : \tau \rangle\rangle \\
      \langle\langle \mathbf{pop} : \tau \rangle\rangle & = & \tau \equiv \langle true, \emptyset \rangle\\
      \langle\langle \mathbf{drop} : \tau \rangle\rangle & = & \tau \equiv \langle true, \emptyset \rangle\\
      \langle\langle \mathbf{dropall} : \tau \rangle\rangle & = & \tau \equiv \langle true, \emptyset \rangle\\
      \langle\langle \mathbf{peek} : \tau \rangle\rangle & = & \tau \equiv \langle true, \emptyset \rangle\\
      \langle\langle \mathbf{peekall} : \tau \rangle\rangle & = & \tau \equiv \langle true, \emptyset \rangle\\
    \end{array}
  \]
  \caption{Extending the constraint generation algorithm}
  \label{fig:constraint-gen}
\end{figure}

We now argue that these extensions interact cleanly with the existing constraint
solver, so that the metatheory of Cardoso et al.~\cite{Cardoso23} carries over to
our setting. The precise property we rely on is their \emph{soundness and
completeness} result for constraint generation: the algorithm
$\langle\langle e : \tau\rangle\rangle$ produces a solvable constraint set exactly
for the expressions that are typable under the declarative type system, and every
solution denotes a valid typing. We must therefore check that neither the new
term formers nor the new side conditions invalidate this correspondence or the
solver that decides it.

The key observation is that the normalized index $\mathbf{norm}^{\ast}(i)$ is a
\emph{static} over-approximation, computed once at constraint-generation time
(Section~\ref{sec:pest-formalization}): it always reduces to a ground value
$\nu \in \{\, n, \mathbf{infty}\,\}$, it never depends on the parsing stack, and it
never contains a unification variable. Consequently the new type formers
$\alpha^{v}$ and $\alpha[v,v]$ appearing in the generation rules above introduce
no index-sort unknowns into the constraint set: the only unification variable they
mention is the fresh body variable $\alpha$, which the solver already handles.
Structurally, $\alpha^{v}$ and $\alpha[v,v]$ behave exactly like the existing type
constructors: their nullability flag and head set are computed field-wise from
$\alpha$ and the ground $\nu$, by the same equations the solver already uses for
sequence, choice and the Kleene star, so the solver needs no new decomposition
rule for them.

The genuinely new material is the two side conditions
$\mathbf{norm}^{\ast}(i) \not\equiv \mathbf{infty}$ and
$\neg(\mathbf{norm}^{\ast}(j) \equiv \mathbf{infty} \land \alpha.\mathrm{null})$. Since
their index arguments are ground, each is a Boolean predicate over the
two-element domain $\{\, n, \mathbf{infty}\,\}$ and the nullability flag, decided
by a constant-time check that is independent of type unification, and can be discharged
either by the SMT backend or by a custom procedure. Such a condition can
never turn an otherwise-solvable constraint set unsolvable, \emph{except} in
precisely the two configurations that the declarative type system also rejects:
an exact repetition with an infinite index, or an unbounded interval over a
nullable body. In other words, the new constraints prune exactly the ill-formed
grammars and no others.

Because the extension (i) adds no new unification variables, (ii) confines index
values to ground positions decided by a separate finite check, and (iii) reuses
the existing nullability and head-set equations, the solver specification
of~\cite{Cardoso23} is unchanged, and their soundness and completeness proof lifts
by the addition of the routine cases for $\alpha^{\nu}$, $\alpha[\nu,\nu]$ and the
ground side conditions. A fully \emph{mechanized} proof of inference soundness and
completeness, complementing the strong-normalization proof of
Section~\ref{sec:typesoundnesslean}, is left as future work. This analysis
leads to our main result:

\begin{Theorem}[Termination of typed grammars]\label{thm:termination}
  Let $G = (V,\Sigma,R, e_s)$ be a well-typed grammar which may use any of Pest
  repetition or stack operators. Then, for every input string $w$ there exists a
  parsing result $r$, either a success or a failure and never a divergent computation, such that $(e_s,w,[])\Rightarrow r$.
\end{Theorem}

The only way for a parsing expression to loop over an input is due to
the presence of (direct or indirect) left-recursive rules or nullable expressions
under a star operator in a grammar~\cite{Ford2004}. Since the type system (and its
inference algorithm) reject any grammar that has these issues, every well-typed
grammar terminates. We do not leave this argument informal: it is proved, for the
full set of Pest operators, as a machine-checked strong-normalization theorem in
Agda (Section~\ref{sec:typesoundnesslean}, Theorem~\lstinline{strong-normalization}).

\section{Mechanising strong normalization in Agda}\label{sec:typesoundnesslean}

The central claim of this paper is that well-typed Pest grammars
\emph{terminate on every input}. In this section we make that claim precise and
prove it, once and for all, in the Agda proof
assistant~\cite{agda}. Rather than reasoning about a fuel-indexed interpreter (whose termination would have to be argued informally, since a fuel-based interpreter always stops when the fuel runs out), we establish a genuine
\emph{strong normalization} result: for a well-typed grammar, the (fuel-free,
relational) parsing judgment relates every input to a result. There is no
gas counter to exhaust and no divergent behaviour to rule out after the fact.

Our proof follows the structure of Tait's reducibility argument for the
strongly-normalising simply-typed $\lambda$-calculus~\cite{plfa22.08}: we define
a \emph{logical relation} on parsing expressions, show that every syntactic
operator preserves it (the \emph{fundamental lemma}), and close the loop for the
mutually-recursive grammar rules by a lexicographic well-founded induction (the
\emph{fixpoint lemma}). The development, which we call \texttt{pest-sn}, extends
a strong-normalization proof for core PEGs to the full set of Pest operators:
the implicit stack, exact repetition and interval repetition.

\subsection{Syntax}\label{sec:agda-syntax}

Parsing expressions are indexed by the number $n$ of nonterminals in the
grammar; a nonterminal is a \lstinline{Fin n}, an index into the rule table. On
top of the seven core PEG constructors we add the six Pest stack operators and
the two repetition operators.

\begin{code}[hide]
module Fig-Syntax where
  open import Data.Char using (Char)
  open import Data.Fin using (Fin)
  open import Data.Nat using (ℕ)
  open import Pest.Index using (Index)
\end{code}
\begin{code}
  data PExp (n : ℕ) : Set where
    empty   : PExp n
    term    : Char → PExp n
    var     : Fin n → PExp n
    seq     : PExp n → PExp n → PExp n
    choice  : PExp n → PExp n → PExp n
    star    : PExp n → PExp n
    not     : PExp n → PExp n
    push    : PExp n → PExp n
    pop peek drop peekall dropall : PExp n
    rep     : Index → PExp n → PExp n
    bound   : Index → Index → PExp n → PExp n
\end{code}

Here \lstinline{push e} parses \lstinline{e} and pushes the consumed prefix onto
the stack; \lstinline{pop}, \lstinline{peek} and \lstinline{peekall} match a
word taken from the stack against the input; \lstinline{drop} and
\lstinline{dropall} discard stack entries; \lstinline{rep i e} is the exact
repetition $e^{i}$; and \lstinline{bound i j e} is the interval repetition
$e[i,j]$. The repetition indexes range over a small arithmetic sublanguage with
literals, $\infty$, the two stack-top queries \textbf{top.tonat} and
\textbf{top.length}, and addition and multiplication; an index is evaluated
against the current stack by a total function
$\mathit{norm} : \mathit{Index} \to \mathit{Stack} \to \mathit{Maybe}\ \mathit{Value}$,
which returns \lstinline{nothing} when the query is undefined (an empty stack, or
a non-numeric top for \textbf{top.tonat}).

\subsection{The type system}\label{sec:agda-typing}

As in Sections~\ref{sec:background} and~\ref{sec:pest-formalization}, a type is a
pair $\langle b, S\rangle$ of a nullability flag and a head set, and a context
$\Gamma$ assigns a type to each nonterminal. The typing judgment
$\mathtt{HasType}\ \Gamma\ e\ \tau$ is the Agda rendering of the rules of
Figures~\ref{fig:typing-repetition} and~\ref{fig:typing-stack}. A grammar is
\emph{well-formed} when there is a context $\Gamma$ that (i) is a fixpoint of the
rule table (every rule body has the type assigned to its nonterminal), and
(ii) is \emph{acyclic}: no nonterminal occurs in its own head set. Acyclicity is
exactly the criterion that rules out (direct and indirect) left recursion, and
together with the non-nullability side condition on \lstinline{star} it is what
makes the head-set cardinality a well-founded measure.

A type is a record of a nullability flag and a head set (a finite subset of the
non-terminals), and a context maps each non-terminal to a type. 
\begin{code}[hide]
module Fig-Typing where
  open import Data.Bool using (Bool; true; false; _∧_; _∨_; if_then_else_)
  open import Data.Fin using (Fin)
  open import Data.Fin.Subset using (Subset; _∉_; _∪_; ⊥; ⁅_⁆)
  open import Data.Nat using (ℕ)
  open import Data.Product using (∃; _,_)
  open import Relation.Binary.PropositionalEquality using (_≡_)
  open import Relation.Nullary using (¬_)
  open import Pest.Syntax
  open import Pest.Index using (Index; infinite; norm-static)
\end{code}
\begin{code}
  record Ty (n : ℕ) : Set where
    constructor mkTy
    field nullable : Bool
          head-set : Subset n
  open Ty

  Ctx : ℕ → Set
  Ctx n = Fin n → Ty n
\end{code}

\clearpage 

The type-forming
operations of Section~\ref{sec:background} become plain functions.

\begin{code}
  _⊗_ _⊕_ : ∀ {n} → Ty n → Ty n → Ty n
  τ₁ ⊗ τ₂ = mkTy (nullable τ₁ ∧ nullable τ₂)
                     (head-set τ₁ ∪ (if nullable τ₁ then head-set τ₂ else ⊥))
  τ₁ ⊕ τ₂ = mkTy (nullable τ₁ ∨ nullable τ₂) (head-set τ₁ ∪ head-set τ₂)

  starTy repTy : ∀ {n} → Ty n → Ty n
  starTy τ = mkTy true (head-set τ)
  repTy τ = mkTy true (head-set τ)
\end{code}
The typing judgment is an inductive family indexed by the expression and its type.
We show the core PEG rules and the representative extensions; the remaining stack
operators (\lstinline{peek}, \lstinline{drop}, \lstinline{peekall},
\lstinline{dropall}) have the same type as \lstinline{pop}, and interval
repetition is typed like exact repetition.
\begin{code}
  data HasType {n} (Γ : Ctx n) : PExp n → Ty n → Set where
    htEmpty  : HasType Γ empty (mkTy true ⊥)
    htTerm   : ∀ {c} → HasType Γ (term c) (mkTy false ⊥)
    htVar    : ∀ {i} → HasType Γ (var i)
                         (mkTy (nullable (Γ i)) (⁅ i ⁆ ∪ head-set (Γ i)))
    htSeq    : ∀ {e₁ e₂ τ₁ τ₂} → HasType Γ e₁ τ₁ → HasType Γ e₂ τ₂
             → HasType Γ (seq e₁ e₂) (τ₁ ⊗ τ₂)
    htChoice : ∀ {e₁ e₂ τ₁ τ₂} → HasType Γ e₁ τ₁ → HasType Γ e₂ τ₂
             → HasType Γ (choice e₁ e₂) (τ₁ ⊕ τ₂)
    htStar   : ∀ {e τ} → HasType Γ e τ → nullable τ ≡ false
             → HasType Γ (star e) (starTy τ)
    htNot    : ∀ {e τ} → HasType Γ e τ → HasType Γ (not e) (mkTy true (head-set τ))
    htPush   : ∀ {e τ} → HasType Γ e τ → HasType Γ (push e) τ
    htPop    : HasType Γ pop (mkTy true ⊥)
    htRep    : ∀ {i e τ} → HasType Γ e τ → ¬ (norm-static i ≡ infinite)
             → HasType Γ (rep i e) (repTy τ)
\end{code}
Finally, a grammar is well-typed when there is a context that types every rule
body at the type it assigns to that non-terminal (the fixpoint condition
\lstinline{hfix}), that is acyclic (\lstinline{hacy}: no non-terminal is in its
own head set, ruling out left recursion), and that types the start expression.

\clearpage

\begin{code}
  record WellTyped {n} (G : Grammar n) : Set where
    field Γ      : Ctx n
          hfix   : ∀ (i : Fin n) → HasType Γ (Grammar.rules G i) (Γ i)
          hacy   : ∀ (i : Fin n) → i ∉ head-set (Γ i)
          hstart : ∃ λ τ → HasType Γ (Grammar.start G) τ
\end{code}
The two repetition operators are typed conservatively as nullable, with the head
set of their body; this discharges their contribution to the progress argument
(their finite iteration count always terminates, whether or not the body is
nullable) while keeping the termination measure untouched. Interval repetition
requires both bounds to be statically finite; the unbounded form $e[i,\infty]$ is
written as the composition of $e^{i}$ and $e^{*}$, where the non-nullability
requirement of the Kleene star applies.

\subsection{Relational semantics}\label{sec:agda-semantics}

We give the semantics as a big-step relation \lstinline{Parses G e s st r},
read: expression \lstinline{e} on input \lstinline{s} with stack \lstinline{st}
yields result \lstinline{r}, where a result either records the consumed prefix,
the remaining suffix and the resulting stack, or is a failure.

\begin{code}[hide]
module Fig-Result where
  open import Pest.Syntax using (Word; Stack)
\end{code}
\begin{code}
  data ParseResult : Set where
    ok   : Word → Word → Stack → ParseResult
    fail : ParseResult
\end{code}

The stack is threaded through the rules but, crucially, plays no part in the
termination argument: the stack operators are all nullable except
\lstinline{push} (which inherits its body's behaviour), so none of them can loop.
Exact and interval repetition are defined by two auxiliary relations,
\lstinline{RepN G e k} (the body run exactly \lstinline{k} times) and
\lstinline{BoundR G e lo d} (the greedy interval $[lo, lo+d]$), which are
mutually inductive with \lstinline{Parses}. Their termination follows by
induction on the finite count and on the interval width, respectively.

\subsection{The logical relation}\label{sec:agda-rel}

The heart of the proof is the reducibility predicate. For a well-formed grammar
\lstinline{G}, context $\Gamma$, type $\tau$ and expression \lstinline{e}, the
relation $\mathcal{R}\ G\ \Gamma\ \tau\ e$ bundles the two invariants we need.

\begin{code}[hide]
module Fig-Rel where
  open import Data.Nat using (ℕ)
  open import Data.List using ([])
  open import Data.Bool using (Bool; false)
  open import Data.Product using (∃)
  open import Relation.Binary.PropositionalEquality using (_≡_; _≢_)
  open import Pest.Syntax
  open import Pest.Typing
  open import Pest.Semantics
\end{code}

\clearpage

\begin{code}
  record 𝓡 {n : ℕ} (G : Grammar n) (Γ : Ctx n)
                    (τ : Ty n) (e : PExp n) : Set where
    field
      terminates : ∀ (s : Word) (st : Stack)
                 → ∃ λ r → Parses G e s st r
      progress   : nullable τ ≡ false
                 → ∀ s st p s' st'
                 → Parses G e s st (ok p s' st') → p ≢ []
\end{code}

The \lstinline{terminates} field is strong normalization for a single
expression: parsing always produces a result. The \lstinline{progress} field is
the PEG analogue of head reduction: a non-nullable expression that succeeds
consumes at least one symbol. Progress is what guarantees that the suffix handed
to the next iteration of a Kleene star (or to the tail of a sequence) is strictly
shorter, so that the well-founded recursion on input length makes progress.

The remaining difficulty is left recursion between grammar rules, handled by a
head-set \emph{rank}: under acyclicity, following a head-dependency edge strictly
decreases the cardinality of the head set. Termination and progress are then
established together by a lexicographic well-founded induction on the pair
$(\text{input length}, \text{head-set cardinality})$; the stack operators are
non-recursive, and the two repetition operators are handled by a standalone
helper that only ever revisits the (structurally smaller) body, so neither
disturbs the measure.

\subsection{The fundamental lemma and strong normalization}\label{sec:agda-thm}

The fundamental lemma states that, in an environment where every grammar rule is
already reducible, every well-typed expression is reducible. The fixpoint lemma
is what supplies that environment. We write $\Gamma \vdash e : \tau$ for the
typing judgment \lstinline{HasType}, $G(i)$ for the body of the rule of the
non-terminal $i$, and $\mathcal{R}_{G,\Gamma}(\tau, e)$ for the reducibility
predicate defined above.

\begin{code}[hide]
open import Pest.Syntax
open import Pest.Index
open import Pest.Typing
open import Pest.Semantics
open import Pest.Rel
open import Pest.Fundamental
open import Pest.Fixpoint
open import Pest.MainTheorem

fundamental-lemma
  : ∀ {n} {G : Grammar n} {Γ : Ctx n}
  → (∀ i → 𝓡 G Γ (Γ i) (Grammar.rules G i))
  → ∀ {e : PExp n} {τ : Ty n}
  → HasType Γ e τ → 𝓡 G Γ τ e
fundamental-lemma = fundamental

fixpoint-lemma
  : ∀ {n} {G : Grammar n} {Γ : Ctx n}
  → (∀ i → i ∉ head-set (Γ i))
  → (∀ i → HasType Γ (Grammar.rules G i) (Γ i))
  → ∀ i → 𝓡 G Γ (Γ i) (Grammar.rules G i)
fixpoint-lemma = fixpoint
\end{code}

\begin{Lemma}[Fundamental lemma]\label{lem:fundamental}
Let $G$ be a grammar with $n$ non-terminals and $\Gamma$ a context such that
$\mathcal{R}_{G,\Gamma}(\Gamma(i), G(i))$ holds for every non-terminal $i$. Then,
for every expression $e$ and type $\tau$, if $\Gamma \vdash e : \tau$ then
$\mathcal{R}_{G,\Gamma}(\tau, e)$.
\end{Lemma}

\begin{proofsketch}
The proof is by structural
induction on the derivation of $\Gamma \vdash e : \tau$, and each case
invokes a closure lemma showing that the corresponding operator preserves the
relation. The base cases exhibit the result directly; sequence and choice combine
the reducibility of the two sub-expressions; the not-predicate inverts the result
of its body; the Kleene star runs a well-founded induction on the length of the
remaining input, using the progress component to ensure that every iteration
consumes at least one symbol, so the input strictly shrinks; the stack operators
are immediate, reducing to their body (\lstinline{push}) or producing a result in
a single step; and the two repetition operators iterate a standalone helper on the
finite count obtained from the stack. The only non-structural information the proof
needs is reducibility for a non-terminal, which is supplied by the environment
hypothesis rather than by recursion through the grammar.
\end{proofsketch}

\begin{Lemma}[Fixpoint lemma]\label{lem:fixpoint}
Let $G$ be a grammar with $n$ non-terminals and $\Gamma$ a context that is
acyclic, that is, $i \notin \mathit{head\text{-}set}(\Gamma(i))$ for every
non-terminal $i$, and that types every rule body at the type it assigns to that
non-terminal, that is, $\Gamma \vdash G(i) : \Gamma(i)$ for every $i$. Then
$\mathcal{R}_{G,\Gamma}(\Gamma(i), G(i))$ holds for every non-terminal $i$.
\end{Lemma}

\begin{proofsketch}
This is the step that cannot be
structural, because expanding a non-terminal replaces it by its rule body, which
is not a subterm of the expression. We proceed by a lexicographic well-founded
induction on the pair of the input length and the head-set cardinality, the
technique introduced in Section~\ref{sec:background}. Expanding a non-terminal
leaves the input unchanged but strictly decreases the head-set rank; this is
legitimate because acyclicity makes the head set of a non-terminal a strict
superset of the head sets reachable from it, a monotonicity property we prove
separately (\lstinline{firstMono}). A Kleene star leaves the type unchanged but
strictly decreases the input length, and every other constructor recurses on a
structurally smaller derivation. Following the development, the argument is split
into a progress part and a termination part, which is what lets Agda's termination
checker accept it.
\end{proofsketch}

Combining the two lemmas yields the main theorem: every well-typed expression of
a well-formed grammar terminates on every input and every initial stack, and in
particular parsing a grammar from its start symbol always terminates. We write
$\mathit{Parses}\ G\ e\ s\ st\ r$ for the big-step judgment of
Section~\ref{sec:agda-semantics}, and $\mathit{start}(G)$ for the start
expression of $G$.

\begin{code}[hide]
strong-normalization
  : ∀ {n} {G : Grammar n} (wt : WellTyped G)
    {e : PExp n} {τ : Ty n}
  → HasType (WellTyped.Γ wt) e τ
  → ∀ (s : Word) (st : Stack) → ∃ λ r → Parses G e s st r
strong-normalization = terminates

start-normalization
  : ∀ {n} {G : Grammar n} → WellTyped G
  → ∀ (s : Word) → ∃ λ r → Parses G (Grammar.start G) s [] r
start-normalization = start-terminates
\end{code}

\begin{Theorem}[Strong normalization]\label{thm:sn-pest}
Let $G$ be a well-typed grammar and $\Gamma$ the context witnessing its
well-typedness. For every expression $e$ and type $\tau$ such that
$\Gamma \vdash e : \tau$, every input $s$ and every stack $st$, there exists a
result $r$ such that $\mathit{Parses}\ G\ e\ s\ st\ r$.
\end{Theorem}

\begin{Theorem}[Termination of parsing]\label{thm:start-pest}
Let $G$ be a well-typed grammar. Then, for every input $s$, there exists a result
$r$ such that $\mathit{Parses}\ G\ \mathit{start}(G)\ s\ [\,]\ r$.
\end{Theorem}

\begin{proofsketch}
Both theorems are immediate from the two lemmas. The fixpoint
lemma produces an environment in which every rule body is reducible; the
fundamental lemma then lifts reducibility to the given well-typed expression; and
the \lstinline{terminates} component of the resulting witness is exactly the
required result. Starting from the grammar's start symbol is the special case in
which the initial stack is empty.
\end{proofsketch}

The two lemmas and the two theorems are stated in the source of this paper as
aliases of the functions \lstinline{fundamental}, \lstinline{fixpoint},
\lstinline{terminates} and \lstinline{start-terminates} of the development, and
are therefore typechecked against it: the mathematical statements above are
transcriptions of the types Agda accepts.

Theorem~\ref{thm:sn-pest} is the mechanised form of the Termination
Theorem~\ref{thm:termination} of Section~\ref{sec:pest-formalization}. Two points
are worth stressing. First, unlike the ``well-typed programs do not get stuck''
statement of a definitional interpreter, it does
not mention fuel: the existential $\exists\,r.\ \mathit{Parses}\ G\ e\ s\ st\ r$
is an unconditional guarantee that a result exists, so termination is a theorem
rather than an informal consequence of a fuel discipline. Second, the result is
about the same big-step judgment used to specify the operators in
Section~\ref{sec:pest-formalization}, so no separate adequacy argument between an
inductive semantics and an interpreter is required.

The development is constructive throughout: it contains no postulates, no uses of
the \lstinline{TERMINATING} pragma and no holes, so Agda's own coverage and
termination checkers vouch for the totality of every function that witnesses the
proof.
\section{Examples}~\label{sec:examples}

We implemented a Racket library, called \emph{typed-peg},
which implements the type inference and the semantics of grammars with the proposed new
operators. In this section, we present a complete example grammar and discuss how
our approach avoid the looping behavior reported by Pest users~\cite{PESTBUG}.

\paragraph{A grammar for the Heartbeat packages} One of the most famous binary data format
bug is the so-called ``heartbleed''. The ASN.1 binary interface uses a ``heartbeat''
package to check if the other remote party is still alive. A well-formed heartbeat package
is formed by the following components:
\begin{itemize}
  \item A type field, which indicates what kind of ``heartbeat'' is being used,
  \item A field for the length $l$, a two byte value,
  \item A challenge $C$ of length $|C| = l$, and
  \item An arbitrary number of padding bytes.
\end{itemize}

\begin{verbatim}
#lang typed-peg
bit <-- '0' / '1'
byte <-- bit^8
type <-- byte
length <-- push(byte^2)
challenge <-- byte^top.tonat
padding <-- byte*
start: type ~ length ~ challenge ~ padding
\end{verbatim}
The previous grammar is a simple encoding of the heartbeat package format.
A key issue regarding implementing parsers for this format is to ensure that
challenge has the size specified by the length field. Using the stack operator
\textbf{push} we insert the parsed length into the execution stack and retrieve
its numeric value using the index expression \textbf{top.tonat}, which parses
the challenge using the exact repetition operator.

\paragraph{A grammar for the PNG format}

A popular image file format is the Portable Network Graphics (PNG), which
is represented by the following pieces of data: 1)
A format signature (8 bytes), which always corresponds to the
following integer values: 137, 80, 78, 71, 13, 10, 26, 10;
2) A 4 bytes natural number $n$, which represents the data length;
3) 4 bytes chunk type value;
4) The image data, a sequence of $n$ bytes and
5) a 4 bytes CRC checksum.
A possible grammar to specify the PNG format is:
\begin{verbatim}
#lang typed-peg
bit <-- '0' / '1'
byte <-- bit^8
fist-four <-- '137' ~ '80' ~ '78' ~ '71'
snd-four <-- '13' ~ '10' ~ '26' ~ '10'
signature <-- first-four ~ snd-four
length <-- push(byte^4)
type <-- byte^4
data <-- byte^top.tonat
crc <-- byte^4
start: signature ~ length ~ type ~ data ~ crc
\end{verbatim}
Notice that the use of an exact repetition operator combined with
index \textbf{top.tonat} allows us to guarantee that the data block
has the correct size of the format length field.

\paragraph{Rejecting reported looping examples} Now, let's turn our
attention to some expressions which Pest users~\cite{PESTBUG} reported as
non terminating.

\begin{verbatim}
forever1 <-- (push('') ~ pop)*
forever2 <-- popall ~ (popall)*
forever3 <-- forever3
\end{verbatim}
In the first two examples, we have nullable expressions under a star
operator. In the first rule, we have a repetition of \textbf{push} operator, which
has an empty string expression as its argument, followed by \textbf{pop}.
A \textbf{push} is nullable only if its argument is nullable. Since it
has an empty string as argument, the expression \verb|push('')|
accepts the empty string. Next, we have a pop operator which is nullable,
by the definition. In this way, the type system rejects \verb|forever1|,
since it has a nullable expression under a star. The second example is also
ill-typed since it has a \textbf{popall}, a nullable operator, as
argument to a PEG repetition.
The third example shows a simple left recursive rule
which is also rejected by type inference algorithm.

\paragraph{A grammar for an indentation sensitive language}

Parsing indentation sensitive languages poses a challenge to language developers.
The next grammar is an example of how stack operators can be used to define
a language which uses different levels of indentation. The idea is to use the
stack to control the indentation level.
\begin{verbatim}
#lang typed-peg
new_line <-- '\n'
char <-- 'a' / ... / 'z' / '0' / ... / '9'
ident_level <-- ' '+ / '\t'+
block_name <-- char+
block <-- block_name ~ new_line ~ new_block*
new_block <-- peekall ~
              push (ident_level) ~
              block ~
              (peekall ~ block)* ~
              drop
start: new_line* ~ block* ~ new_line*
\end{verbatim}
At each block, we start by matching the current block indentation
using the operator \textbf{peekall}. Next, we push the possible
next indentation level and process a possibly empty set of blocks.
We end a block by removing its indentation from the parser stack
using \textbf{drop}.

\paragraph{Empirical evaluation} A natural question is whether the type system is
too restrictive on grammars written in practice. To answer it we ran the
inference algorithm on a corpus of \emph{real} Pest grammars harvested from
projects listed in the \texttt{pest-parser/awesome-pest} collection: the pest
project's own example, test and reference grammars (JSON, TOML, SQL, HTTP, a
calculator, indentation-sensitive lists, and others), together with grammars from
several real world Rust projects.
We record for each grammar its source repository, the path of the original
\texttt{.pest} file and the commit it was taken from.

Pest's surface syntax is richer than the fragment our prototype implements, so
each \texttt{.pest} file is translated to the prototype syntax.
The translation is \emph{sound for the type system} because the
analysis depends only on nullability and head-sets: every single-character
matcher (a string or character literal, a character range, an \texttt{ASCII\_*}
built-in, or \texttt{ANY}) has the type $\langle false, \emptyset\rangle$, so all
of them are abstracted to a single terminal, while zero-width anchors become
$\epsilon$ and the constructs that actually drive the analysis (non-terminal
references, sequence, choice, repetition, predicates and the stack operators) are
preserved. Grammars that are fragments of a multi-file grammar, and hence contain
undefined references, or that use a construct outside our fragment (for instance
the Pratt-parser rules), are reported as \emph{not-expressible} rather than
silently dropped.

Of 32 grammar files, 25 were analyzable and \emph{all 25 were accepted} as
well-formed; \emph{none was rejected}. The remaining 7 are not-expressible,
because they reference rules defined in a companion grammar file or in Rust code.
In other words, on the real grammars that the prototype can analyze, the type
system is not restrictive at all: every one of them is well-typed. This does not
mean the type system is vacuous. Its ability to reject looping grammars is
exercised by a second, deliberately adversarial corpus of grammars that are
designed to loop, over which the algorithm rejects exactly the patterns it is
meant to exclude (a nullable body under a Kleene star, direct and indirect left
recursion, and an exact repetition with an infinite index), with no false
positives. Together the two corpora indicate that the type system flags the
genuine sources of non-termination while remaining permissive on real-world
grammars.

% While preparing this evaluation we hardened the prototype. We fixed three
% implementation defects, namely a mismatched constructor name that rejected every
% use of the \textbf{dropall} operator and two parser bugs in the handling of
% character classes and of the optional operator, and we re-enabled the
% one-or-more operator. We also replaced the fragile textual parsing of the SMT
% solver output by a direct solver interface built on the Rosette solver-aided
% language~\cite{rosette}. The previous implementation serialised the constraints
% to an SMT-LIB script, invoked an external \texttt{z3} executable, and parsed its
% textual model with a hand-written parser, which breaks on Z3 versions newer than
% the one it was written for; on five grammars of the corpus the original parser
% failed on the solver's output while the Rosette interface returned a verdict. The
% Rosette interface removes the dependency on an external executable and the need to
% parse solver output, and it agrees with the original solver on every grammar that
% both can decide. The corpus, the provenance manifest and the scripts that
% reproduce these numbers are available in the artifact.

\paragraph{Discussion}
Our work uses the framework of operational semantics and type
theory to provide a solid foundation to Pest stack and repetition
operators. The main issue pointed by Pest users~\cite{PESTBUG} is failure of
detecting non termination of grammars. The proposed type inference
algorithm can avoid this problem by ruling out any grammar which
has direct or indirect left-recursive rules and nullable expressions
as a parameter to the Kleene star operator.

During our formalization, we noticed that the inclusion
of indexes that allow the inspection of the topmost stack item
would increase the tool expressivity, since they allow concise
specifications for some context-sensitive formats. We must, however, be honest
about a limitation of the present type system: while it guarantees termination,
it does \emph{not} rule out all run-time errors of the two new indexes. Evaluating
$\mathbf{top.tonat}$ or $\mathbf{top.length}$ fails when the stack is empty, and
$\mathbf{top.tonat}$ additionally fails when the top element is not a string of
digits. In our mechanized semantics these situations do not cause divergence, since the parser simply reports a failure (the \lstinline{norm} function returns
\lstinline{nothing} and the repetition rule produces a \lstinline{fail} result,
so strong normalization is preserved), but a stronger type system could reject
such grammars statically. Since these two indexes are proposed by us and are not
yet part of Pest, ensuring their safety is a natural next step.

We sketch how the type system could be extended to that end. The essential
missing information is a \emph{coarse abstraction of the stack} that each
expression operates on. Concretely, one would enrich the type of an expression
with an abstract stack: a (bounded) list of abstract words, where an abstract
word records whether it is guaranteed to be non-empty, whether it is guaranteed to
denote a number (for $\mathbf{top.tonat}$), and, when useful, a bound on its
length. The stack operators would then act on this abstraction ($\mathbf{push}(e)$
pushes an abstract word derived from the type of $e$, while $\mathbf{pop}$ and
$\mathbf{drop}$ pop it), and the typing rule for $\mathbf{top.tonat}$ (resp.
$\mathbf{top.length}$) would require the abstract stack to be non-empty with a
numeric (resp. any) top. This is an effect-like discipline tracking stack shape
rather than stack contents; a deliberately coarse abstraction (for instance,
tracking only the top-of-stack word, or a fixed-depth prefix) keeps it decidable
and composable with the head-set inference of Section~\ref{sec:pest-formalization},
at the cost of rejecting some safe but data-dependent grammars. Designing and
mechanizing such a stack-aware type system is left for future work.

A further limitation of our current prototype is that it does not handle different
input encodings: the semantics of $\mathbf{top.tonat}$ depends on the current
input encoding to convert the topmost stack value to a natural number. The proper
handling of multiple encodings is also left for future work.

\section{Related work}~\label{sec:related}

The works of Ribeiro et al. ~\cite{Ribeiro19} and Cardoso et. al.~\cite{Cardoso23}
proposed a type system and a type inference algorithm for PEGs that is sound with
respect to Ford's well-formedness predicate. Authors argue that the use
of a type system provide a more predictable and compositional behavior.
Another work that uses types to ensure parsing termination was proposed by
Krishnaswami et. al.~\cite{Krishnaswami19}. The authors define a type system
for $\mu$-regular expressions, an algebraic presentation of context-free
languages, and proved that their type system only accepts LL(1) grammars.
Our work extends both the type system for PEGs and its inference algorithm to
Pest operators while ensuring that no well-typed grammar
will loop over inputs.

The development of formalisms and tools for correct parsing of complex data
formats is subject of some recent works.  Zhang et. al.~\cite{Zhang23} proposed
interval parsing grammars (IPG) as formalism for expressing some context-sensitive
formats present in binary data, like video and image formats. An IPG attaches to
every non terminal/terminal an interval, which specifies the range of input consumed.
By connecting intervals and attributes, the context-sensitive patterns in file formats
can be well handled. Unlike \emph{typed-peg}, IPGs allow the use semantic actions to
perform arbitrary computations over attributes. While such approach provide extra
power, in our view, it clutters the grammar since the parser specification is not formed only by
grammar rules but also program pieces to handle parsed data.

Another work that proposes a formalism to deal with data formats was developed by Lucks
et. al.~\cite{Lucks17}. Authors devised a new class of languages, named
``calc-regular languages'', which enjoy the property of prefix-freeness that is
present in several binary data formats. In order to provide concise specification
of calc-regular languages, author devised \emph{calc-regular expressions}, an
extension of regular expressions. The parsing of such languages can be made by
finite state machines with accumulators, which are proved equivalent to calc-regular
expressions. Zephyr et. al~\cite{Zephyr21} proposed an extension to PEGs, based on
calc-regular expressions, named calc-PEGs which provide a specific length
operator to PEGs. Authors discuss a $O(n^2)$ algorithm to parse calc-PEGs and
present a tool, called pegmatite, that produces parsers from calc-PEGs specifications.
The length operator of calc-PEGs is close in spirit to our stack-top indexes: both
let a grammar constrain the number of symbols matched by a sub-expression using a
value read from the input. The difference is one of mechanism and scope. Their
length operator binds a numeric measure of a matched fragment and reuses it as a
repetition count within the same expression, whereas our \textbf{top.tonat} and
\textbf{top.length} read a value that was explicitly \emph{pushed} onto the
implicit stack (its numeric interpretation, respectively its length), decoupling
where a length is measured from where it is later used and composing with the
other stack operators. Unlike our work, the authors offer no guarantee of
termination, nor do they provide a detailed specification of calc-PEGs.

The use of intrisincally-typed interpreters in semantics are the subject of
some recent research works~\cite{Feitosa19, Feitosa22, Thiemann23, Rouvoet20,
Rest22}.
Feitosa and Ribeiro~\cite{Feitosa22} present a mechanically verified
implementation of the list-machine benchmark using the dependently typed language
Agda. Authors provide provably correct algorithms for subtyping and least common
supertype calculation, as well as a type checker that produces intrinsically-typed
program representations. They show that this approach reduces code size significantly
(using only 14\% of the lines of code compared to a Twelf solution and 47\% compared
to Coq) while still covering all 14 benchmark tasks. The paper also extends the
formalization to support indirect jumps (version 2.0 of the benchmark), demonstrating
how the intrinsic approach handles more complex control flow features like continuation
types and mutually inductive relations. The formalization of Featherweight Java (FJ)
using intrisincally-typed syntax was the subject of~\cite{Feitosa19}.
Authors used a novel encoding that integrates FJ's nominal type system, class tables,
and inheritance into an intrinsically-typed syntax, where only well-typed expressions
are representable. By using de Bruijn indices for names and splitting class definitions
into signatures and implementations, the authors handle complex binding structures and
mutual dependencies elegantly. They further implement a fuel-based definitional
interpreter that leverages the type information to guarantee progress and preservation
theorems without separate proofs.
Thiemann's work~\cite{Thiemann23} uses a definitional interpreter to define a callback-based approach for
implementing binary session types in dependently typed functional languages like Agda,
eliminating the need for linear types in the host language. By encapsulating all
communication logic within an interpreter that forms the trusted computing base,
the approach guarantees protocol fidelity and linearity by construction. The work extends
to advanced session type features including branching, recursion, multichannel sessions,
higher-order sessions, and context-free sessions, while also introducing novel constructs
such as dynamic selection and an UNROLL command for recursive client protocols.
Another work which uses a definitional interpreter-based semantics for a linear and session-typed
language is~\cite{Rouvoet20}. The authors use a syntax encoding which takes inspiration from separation
logic, they develop proof-relevant separation algebras and composable
monadic abstractions (including linear Reader and State monads, and a free
monad for concurrency) that encapsulate the complexities of resource
separation and linearity. These abstractions enable the concise specification of
interpreters for increasingly complex languages: a linear lambda calculus,
one with strong references, and finally a session-typed language with
concurrency and asynchronous communication. The work demonstrates that
with the right abstractions, intrinsically-typed interpreters for sophisticated linear
languages can retain the clarity and simplicity of those for simple languages, while
providing executable and type-safe semantics by construction.

\section{Conclusion}~\label{sec:conclusion}

In this work we proposed a formalization of an extended version of Pest style PEG
grammars. We provide a formalization of stack manipulation and repetition operators
for Pest grammars and presented a mechanized relational semantics together with a
machine-checked \emph{strong normalization} result in Agda, which guarantees
termination of parsing for all well-typed grammars without appealing to a fuel
counter.
Besides a formal definition of the operator semantics, we
extended a typing inference algorithm for PEGs~\cite{Cardoso23} which guarantee
termination of well-typed grammars to handle Pest operators.

As future work, we intend to continue development of the \emph{typed-peg} Racket
library and to integrate the ideas presented in this work into future versions of
the Pest tool. On the formal side, three directions stand out: a stack-aware type
system that also rules out the run-time errors of $\mathbf{top.tonat}$ and
$\mathbf{top.length}$ by tracking a coarse abstraction of the stack; a mechanized
proof of soundness and completeness of the inference algorithm, complementing the
strong-normalization proof of this paper; and support for multiple input
encodings in the semantics of $\mathbf{top.tonat}$.

The complete formalization source code is available on-line~\cite{pegrepository}.

\section*{ACKNOWLEDGMENTS}
This work is supported by the FAPEMIG - Fundação de Amparo à Pesquisa de Minas Gerais -
under grant number APQ-01683-21.


%% bibitems, please use
%%
\bibliographystyle{elsarticle-num}
\bibliography{references}

%% else use the following coding to input the bibitems directly in the
%% TeX file.

%% Refer following link for more details about bibliography and citations.
%% https://en.wikibooks.org/wiki/LaTeX/Bibliography_Management

\end{document}

\endinput
%%
%% End of file `elsarticle-template-num.tex'.
